% Mechanically extracted after \mainmatter from the frozen standalone master.

\markboth{第 XIV 讲. 正则元素（续）}{第 XIV 讲. 正则元素（续）}
% BEGIN inputs/XIV/tex/components/00_expose_XIV_section1_opening_en.tex
% Authority: Expo14.pdf, local pp. 1--3 / combined pp. 835--837.
% Continuous English currently runs through the statement of Lemma 2.1.

\SGARefTarget{sga3:xiv:expose}

\setcounter{footnote}{0}
\footnotetext{2008 年 12 月 1 日 xy 版.}

\section{光滑代数群的 Cartan 子群与极大环面的构造}
\SGARefTarget{sga3:xiv:section:1}

\paragraph{\textbf{定理 1.1.}}
\SGARefTarget{sga3:xiv:1.1}
\emph{设 \(G\) 是域 \(k\) 上的光滑代数群. 则 \(G\) 容许极大环面 \(T\), 因而容许 Cartan 子群 \(C=\Centr_G(T)\).}

由 \SGARefLink{sga3:xii:3.2}{XII~3.2}, 寻找 \(G\) 的极大环面 \(T\) 等价于寻找 \(G\) 的 Cartan 子群 \(C\). 此外, 由于 \(G\) 的极大环面就是 \(G^0\) 的极大环面, 可假定 \(G\) 连通. 区分两种情形.

\medskip
\noindent\textup{1\textdegree)} \emph{域 \(k\) 有限.}
设 \(\mathcal T\) 为 \(G\) 的极大环面概形（\SGARefLink{sga3:xii:1.10}{XII~1.10}）; 它在 \(k\) 上光滑. 群 \(G\) 通过内自同构作用于 \(\mathcal T\), 共轭定理（\SGARefLink{sga3:xii:6.6:item:a}{XII~6.6(a)}）说明, \(\mathcal T_{\bar k}\) 的任意两个 \(\bar k\)-有理点在 \(G_{\bar k}(\bar k)\) 的作用下共轭. 由于 \(\mathcal T\) 在 \(k\) 上光滑, \(\mathcal T_{\bar k}\) 在 \(\bar k\) 上光滑; 从而 \(\mathcal T_{\bar k}\) 同构于
\[
  G_{\bar k}/N,
\]
其中 \(N\) 为 \(\mathcal T_{\bar k}(\bar k)\) 的某个元素的稳定子, 即 \(G_{\bar k}\) 的一个极大环面 \(\bar T\) 的正规化子. 因而 \(\mathcal T\) 是 \(G\) 作用下的齐性空间. Lang 的一个熟知定理（\emph{Amer.\ J.\ Math.} \textbf{78} (1956), 555--563）说明, 有限域 \(k\) 上光滑连通代数群的每个齐性空间都有有理点. 特别地, \(\mathcal T\) 有有理点; 等价地, \(G\) 有极大环面 \(T\).

\medskip
\noindent\textup{2\textdegree)} \emph{域 \(k\) 无限.}
将用到下述引理.

\paragraph{\textbf{引理 1.2.}}
\SGARefTarget{sga3:xiv:1.2}
\emph{设 \(G\) 是域 \(k\) 上的光滑代数群. 则 \(G\) 容许一个 \textup{(C)} 型子群\textup{（\SGARefLink{sga3:xiii:definition:6:2}{XIII~6.2}）}.}

由 \SGARefLink{sga3:xiii:corollary:6:3}{XIII~6.3}, 这等价于说 \(\mathfrak g\) 包含一个 Cartan 子代数 \(\mathfrak d\). 当 \(k\) 无限时, 这是直接的: \(\mathfrak g\) 包含一个正则元素 \(a\), 只须取
\[
  \mathfrak d=\Nil(a,\mathfrak g).
\]
\(k\) 有限的情形与上面 \textup{1\textdegree)} 的证明完全相同, 但须先构造 \(\mathfrak g\) 的 Cartan 子代数概形 \(\mathcal D\), 并证明该概形在 \(k\) 上光滑; 这将在下面 \SGARefLink{sga3:xiv:2.16}{2.16} 中完成.

所以, 为在当前所考虑的情形 \textup{2\textdegree)} 中证明 \SGARefLink{sga3:xiv:1.1}{1.1}, 只须知道当 \(k\) 无限时 \SGARefLink{sga3:xiv:1.2}{1.2} 成立. 还记下下述结果供以后使用.

\paragraph{\textbf{引理 1.3.}}
\SGARefTarget{sga3:xiv:1.3}
\emph{设 \(G\) 为光滑连通仿射代数群, 其约化中心\textup{（\SGARefLink{sga3:xii:4.1}{XII~4.1} 和 \SGARefLink{sga3:xiv:4.4}{4.4}）}退化为单位群. 则 \(G\) 幂零当且仅当其 Lie 代数 \(\mathfrak g\) 幂零.}

这包含在 \SGARefLink{sga3:xii:4.9}{XII~4.9} 之中.

现在可以给出构造 \(G\) 的 Cartan 子群的一种步骤. 当 \(k\) 有限时, 只要接受该情形下的\SGARefLink{sga3:xiv:1.2}{引理~1.2}, 同一步骤亦适用. 先假定 \(G\) 仿射. 对 \(n=\dim G\) 归纳, 当 \(n=0\) 时断言显然成立. 因而假定 \(n>0\), 并且已在维数 \(n'<n\) 时证明断言. 令 \(Z\) 为 \(G\) 的约化中心, 令
\[
  u\colon G\longrightarrow G'=G/Z
\]
为典范同态. 由 \SGARefLink{sga3:xii:theorem:4-7:item:c}{XII~4.7(c)}, 对应
\[
  C'\longmapsto C=u^{-1}(C')
\]
给出从 \(G'\) 的 Cartan 子群到 \(G\) 的 Cartan 子群的双射. 以 \(G'\) 代替 \(G\), 因而可假定 \(G\) 的约化中心 \(Z\) 退化为单位元\textup{（由 \SGARefLink{sga3:xii:theorem:4-7:item:b}{XII~4.7(b)}, \(G'\) 的约化中心亦有此性质）}.

由\SGARefLink{sga3:xiv:1.2}{引理~1.2},\footnote{底本此处印为“\SGARefLink{sga3:xiv:1.3}{1.3}”; 所引的存在性结果是\SGARefLink{sga3:xiv:1.2}{引理~1.2}.} \(G\) 容许一个 \textup{(C)} 型子群 \(D\). 在 \(k\) 的代数闭包 \(\bar k\) 上, 群 \(D_{\bar k}\) 包含 \(G_{\bar k}\) 的一个 Cartan 子群\textup{（\SGARefLink{sga3:xiii:remarks:6:6:part:b}{XIII~6.6(b)}）}; 所以 \(D\) 的每个 Cartan 子群都是 \(G\) 的 Cartan 子群\textup{（\SGARefLink{sga3:xiii:corollary:2:8:item:a}{XIII~2.8(a)}）}. 因而只须寻找 \(D\) 的一个 Cartan 子群.

若 \(\dim D=\dim G\), 即 \(D=G\), 则 \(G\) 的 Lie 代数是自身的 Cartan 子代数, 因而幂零. 此时\SGARefLink{sga3:xiv:1.3}{引理~1.3} 表明 \(G\) 幂零, 所以 \(G\) 是自身的 Cartan 子群. 若 \(\dim D<\dim G\), 则归纳假设给出 \(D\) 的一个 Cartan 子群, 从而也是 \(G\) 的 Cartan 子群. 这就证明了当 \(G\) 仿射时断言成立.

一般地, 令 \(Z\) 为 \(G\) 的中心. 则 \(G/Z=G'\) 仿射\textup{（\SGARefLink{sga3:xii:6.1}{XII~6.1}）}, 而 \(G'\) 的每个 Cartan 子群 \(C'\) 在 \(G\) 中的逆像都是 \(G\) 的 Cartan 子群\textup{（\SGARefLink{sga3:xii:6.6:item:e}{XII~6.6(e)}）}. 所以只须寻找光滑仿射代数群 \(G'\) 的一个 Cartan 子群, 这是已处理的情形.

\paragraph{\textbf{推论 1.4.}}
\SGARefTarget{sga3:xiv:1.4}
\emph{设 \(G\) 是 Artin 概形 \(S\) 上的光滑有限型群概形. 则 \(G\) 容许一个极大环面 \(T\), 因而容许 Cartan 子群 \(C=\Centr_G(T)\). \(G\) 的每个环面都包含在一个极大环面中.}

事实上, 可假定 \(S\) 局部, 其剩余域为 \(k\). 由 \SGARefLink{sga3:xiv:1.1}{1.1}, 群
\[
  G_0=G\times_S\operatorname{Spec}(k)
\]
容许极大环面 \(T_0\). 由 \SGARefLink{sga3:ix:theorem:3:6}{IX~3.6}~bis 和 \SGARefLink{sga3:x:corollary:2:3}{X~2.3}, \(T_0\) 来自 \(G\) 的一个环面 \(T\), 它显然极大. 将前述结果应用于给定环面的中心化子即得最后的断言, 该中心化子在 \(k\) 上光滑\textup{（\SGARefLink{sga3:xi:corollary:2:4}{XI~2.4}）}.

\paragraph{\textbf{注 1.5.}}
\SGARefTarget{sga3:xiv:1.5}
\textup{(a)} 下面将在 \SGARefLink{sga3:xiv:3.20}{3.20}、\SGARefLink{sga3:xiv:3.21}{3.21} 以及第 XV 讲给出 \SGARefLink{sga3:xiv:1.4}{1.4} 的若干变式, 其中不假定 \(S\) 为 Artin 概形.

\medskip
\noindent\textup{\SGAInlineRefTarget{sga3:xiv:remark:1-5:item:b}(b)} 我不知道, 域 \(k\) 上每个未必光滑的代数群 \(G\) 是否都容许极大环面. 问题只在特征 \(p>0\) 时出现. 将 \SGARefLink{sga3:xiv:1.1}{1.1} 应用于形如 \(G'=G/I\) 的光滑商群, 其中 \(I\) 是 \(G\) 的适当的根式子群\textup{（例如 Frobenius 同态的一个充分高次幂的核）}, 再在 \(G\) 中取 \(G'\) 的极大环面的逆像, 就将问题归结到 \((G_{\bar k})_{\mathrm{red}}\) 为环面的情形; 此处一如往常, \(\bar k\) 表示 \(k\) 的代数闭包. 容易看出, 当 \(G\) 交换, 或更一般地幂零时, 答案肯定: 在此情形中 \(G\) 有唯一的极大环面, 例如可将 \(G_{\bar k}\) 的极大环面下降以构造它.\begingroup
\renewcommand{\thefootnote}{*}\footnote{M. Raynaud 对此处提出的问题给出了否定答案; 参见第 XVII 讲, \SGARefLink{sga3:xvii:subsection:5-9:item:c}{例~5.9(c)}.}
\endgroup

\medskip
\noindent\textup{(c)} 当 \(G\) 仿射且 \(k\) 完美或 \(G\) 可解时, Rosenlicht 的一个定理已给出 \SGARefLink{sga3:xiv:1.1}{1.1}; 他的证明与这里所给的证明很不相同.

\medskip
\noindent\textup{(d)} 当 \(k\) 无限时, \SGARefLink{sga3:xiv:1.1}{1.1} 可由下面 \SGARefLink{sga3:xiv:6.1}{6.1} 中所证明的更精确的结果推出: \(G\) 的极大环面概形 \(\mathcal T\) 为有理簇. 其方法本质上是将 \SGARefLink{sga3:xiv:1.1}{1.1} 的证明与 \(\mathfrak g\) 的 Cartan 子代数概形 \(\mathcal D\) 的结构的显式描述相结合. 为获得所需结果, 须先将 XIII~4--6 的某些结果推广到任意底预概形\textup{（这是下两节的目的）}, 然后利用 \(G\) 的每个 Cartan 子群均包含在 \(G\) 的唯一一个 \textup{(C)} 型子群中这一事实\textup{（第~4、5 节）}, 进一步细化前述证明 \SGARefLink{sga3:xiv:1.1}{1.1} 的构造.

\section{任意预概形上的 Lie 代数: 正则截面与 Cartan 子代数}
\SGARefTarget{sga3:xiv:section:2}

\paragraph{\textbf{引理 2.1.}}
\SGARefTarget{sga3:xiv:2.1}
\emph{设 \(A\) 是环, \(\mathfrak d\) 是 \(A\) 上的 Lie 代数, 对每个 \(s\in\operatorname{Spec}(A)\), 令
\[
  \mathfrak d(s)=\mathfrak d\otimes_A k(s)
\]
为剩余域 \(k(s)\) 上的相应 Lie 代数. 假定 \(A\)-模 \(\mathfrak d\) 为有限表示. 下列条件等价:
\begin{enumerate}
\item[\textup{(i)}] 对每个 \(s\in\operatorname{Spec}(A)\), Lie 代数 \(\mathfrak d(s)\) 幂零;
\item[\textup{(ii)}] 对每个 \(x\in\mathfrak d\), 自同态 \(\operatorname{ad}(x)\) 幂零;
\item[\textup{(iii)}] 存在整数 \(N\geq 0\), 使得对 \(\mathfrak d\) 中的每个元素序列 \(x_1,\ldots,x_N\), 都有
  \[
    \operatorname{ad}(x_1)\operatorname{ad}(x_2)\cdots
    \operatorname{ad}(x_N)=0.
  \]
\end{enumerate}}

% END inputs/XIV/tex/components/00_expose_XIV_section1_opening_en.tex

% BEGIN inputs/XIV/tex/components/01_expose_XIV_lemma21_proof_through_definition25_en.tex
% Authority: Expo14.pdf, local pp. 3--4 / combined pp. 837--838.
% Continuous English ends immediately before Proposition 2.6.

当 \(A\) 是域时, \textup{(i)} 与 \textup{(iii)} 的等价性就是幂零的定义, 而 \textup{(ii)} 与 \textup{(iii)} 的等价性是 Engel 定理的熟知推论\textup{（Bourbaki, \emph{Groupes et algèbres de Lie}, 第 I 章, \S4, 第~2 节）}. 一般地, \textup{(iii)} 显然蕴含 \textup{(ii)}, 而由前述结果以及 \textup{(ii)} 在取商和局部化下保持不变这一事实, \textup{(ii)} 蕴含 \textup{(i)}. 只须证明 \textup{(i)} 蕴含 \textup{(iii)}.

先假定 \(A\) 是局部 Artin 环, 极大理想为 \(\mathfrak m\). 取 \(n>0\) 使得 \(\mathfrak m^n=0\), 并选取 \(N\) 使得条件 \textup{(iii)} 对
\[
  \mathfrak d(s)=\mathfrak d\otimes_A A/\mathfrak m.
\]
成立. 取 \(N'=nN\), 立即可见 \(N'\) 对 \(\mathfrak d\) 满足 \textup{(iii)}.

若 \(A\) 为 Noether 环, 则存在有限个点 \(s_i\in\operatorname{Spec}(A)\) 及定义理想 \(\mathfrak q_i\subset A_{s_i}\), 使得自然映射
\[
  \mathfrak d\longrightarrow
  \prod_i\bigl(\mathfrak d\otimes_A A_i\bigr),
  \qquad A_i=A_{s_i}/\mathfrak q_i,
\]
为单射. 由前一种情形, 对每个 \(i\), 都存在对 Lie 代数 \(\mathfrak d\otimes_A A_i\) 满足 \textup{(iii)} 的整数 \(N_i\). 取 \(N\) 为各 \(N_i\) 中最大的一个, 即得对 \(\mathfrak d\) 的条件 \textup{(iii)}. 最后, 利用 EGA~IV~8 中描述的通常步骤, 将一般情形归结到 Noether 情形.

\paragraph{\textbf{定义 2.2.}}
\SGARefTarget{sga3:xiv:2.2}
设 \(S\) 是预概形, \(\mathfrak d\) 是 \(S\) 上的拟凝聚 Lie 代数, 其底层模为有限表示. 若对每个 \(s\in S\), \(k(s)\) 上的 Lie 代数 \(\mathfrak d(s)\) 幂零, 则称 \(\mathfrak d\) \emph{局部幂零}. 若 \(\mathfrak d\) 局部自由, 且在 \(S\) 的每个使其秩为常数 \(r\) 的开子集 \(U\) 上, 其 Killing 多项式退化为
\[
  P_{\mathfrak d}(t)=t^r.
\]

则称它\emph{严格局部幂零}.
当然, 当 \(\mathfrak d\) 是 \(S\) 上的 Lie 代数且作为模局部自由时, 它的 Killing 多项式定义为多项式
\[
  P_{\mathfrak d}\in A[t],
  \qquad
  A=\Gamma\!\left(\operatorname{Sym}_{\mathcal O_S}(\mathfrak d)\right)
   \simeq \Gamma\!\left(W(\mathfrak d)\right),
\]
其中 \(A\) 是由 \(\mathfrak d\) 定义的向量丛 \(W(\mathfrak d)\) 的结构层的截面环.

\SGARefLink{sga3:xiv:2.2}{2.2} 中引入的两个概念显然在基变换下稳定, 并对 fpqc 拓扑是局部的. 记下:

\paragraph{\textbf{命题 2.3.}}
\SGARefTarget{sga3:xiv:2.3}
\emph{若 \(\mathfrak d\) 严格局部幂零, 则它局部幂零. 当 \(\mathfrak d\) 局部自由且 \(S\) 约化时, 逆命题成立.}

证明是直接的.

\paragraph{\textbf{定义 2.4.}}
\SGARefTarget{sga3:xiv:2.4}
设 \(S\) 是预概形, \(\mathfrak g\) 是 \(S\) 上的 Lie 代数, 其底层模为局部自由有限型模, \(\mathfrak d\) 是 \(\mathfrak g\) 的 Lie 子代数. 若 \(\mathfrak d\) 满足下列条件, 则称它为 \emph{Cartan 子代数}:
\begin{enumerate}
\item[\textup{\SGAInlineRefTarget{sga3:xiv:2.4:item:i}(i)}] \(\mathfrak d\) 局部为直和因子子模\textup{（因而本身亦为局部自由有限型模）};
\item[\textup{(ii)}] 对每个 \(s\in S\), Lie 代数 \(\mathfrak d(s)\) 是 \(\mathfrak g(s)\) 的 Cartan 子代数.
\end{enumerate}

\paragraph{\textbf{定义 2.5.}}
\SGARefTarget{sga3:xiv:2.5}
设 \(S\) 与 \(\mathfrak g\) 如 \SGARefLink{sga3:xiv:2.4}{2.4}. 若 \(\mathfrak g\) 的截面 \(a\) 满足: 对每个 \(s\in S\), 元素 \(a(s)\in\mathfrak g(s)\) 在 \(k(s)\) 上的 Lie 代数 \(\mathfrak g(s)\) 中正则, 则称其为\emph{拟正则截面}. 若 \(a\) 拟正则且包含在 \(\mathfrak g\) 的某个 Cartan 子代数中, 则称其为\emph{正则截面}.

\SGARefLink{sga3:xiv:2.4}{2.4} 与 \SGARefLink{sga3:xiv:2.5}{2.5} 的概念同样在基变换下稳定, 并对 fpqc 拓扑是局部的. 只有最后一个断言, 且仅就正则截面的概念而言, 需要证明. 它来自包含给定正则截面的 Cartan 子代数的唯一确定性. 更精确地说:

% END inputs/XIV/tex/components/01_expose_XIV_lemma21_proof_through_definition25_en.tex

% BEGIN inputs/XIV/tex/components/02_expose_XIV_proposition26_corollary27_en.tex
% Authority: Expo14.pdf, local p. 5 / combined p. 839.
% Continuous English ends immediately before Corollary 2.8.

\paragraph{\textbf{命题 2.6.}}
\SGARefTarget{sga3:xiv:2.6}
\emph{设 \(S\)、\(\mathfrak g\) 如 \SGARefLink{sga3:xiv:2.4}{2.4}, \(a\) 为 \(\mathfrak g\) 的拟正则截面. \(\mathfrak g\) 中至多有一个包含 \(a\) 的 Cartan 子代数. 这样的 Cartan 子代数存在——即 \(a\) 是正则截面——当且仅当 \(a\) 满足如下条件:
\[
  \mathfrak d=\Nil(a,\mathfrak g)
\]
局部为 \(\mathfrak g\) 的直和因子子模, 且 \(\operatorname{ad}(a)\) 诱导 \(\mathfrak g/\mathfrak d\) 的自同构. 此时 \(\mathfrak d\) 就是 \(\mathfrak g\) 中唯一一个包含 \(a\) 的 Cartan 子代数.}

先假定 \(a\) 包含在 \(\mathfrak g\) 的某个 Cartan 子代数 \(\mathfrak d\) 中. 则 \(\operatorname{ad}(a)_{\mathfrak g/\mathfrak d}\) 在每个纤维上均为双射. 由于 \(\mathfrak g/\mathfrak d\) 为局部自由有限型模, 这个自同态是 \(\mathfrak g/\mathfrak d\) 的自同构. 另一方面, 由 \SGARefLink{sga3:xiv:2.1}{2.1}, 自同态 \(\operatorname{ad}(a)_{\mathfrak d}\) 局部幂零. 从而
\[
  \mathfrak d=\Nil(a,\mathfrak g),
\]
这同时证明了 \(\mathfrak d\) 的唯一性和所述正则性判据的必要性.

对于充分性, 注意关于 \(a\) 的假设蕴含 \(\Nil(a,\mathfrak g)\) 的形成与底扩张可交换, 特别地与取纤维可交换. 因而 \(\mathfrak d=\Nil(a,\mathfrak g)\) 的各纤维 \(\mathfrak d(s)\) 都是 \(\mathfrak g(s)\) 的 Cartan 子代数. 此外, 由 \SGARefLink{sga3:xiii:proposition:4:1}{XIII~4.1}, \(\mathfrak d\) 是 \(\mathfrak g\) 的 Lie 子代数. 因此它是 Cartan 子代数.

\paragraph{\textbf{推论 2.7.}}
\SGARefTarget{sga3:xiv:2.7}
\emph{在 \SGARefLink{sga3:xiv:2.4}{2.4} 的假设下, 设 \(\mathfrak d\) 为 \(\mathfrak g\) 的 Cartan 子代数, \(a\) 为 \(\mathfrak d\) 的一个在 \(\mathfrak g\) 中正则的截面, \(u\) 为 \(\mathfrak g\) 的自同构. 则 \(u\) 保持 \(\mathfrak d\) 当且仅当 \(u(a)\) 是 \(\mathfrak d\) 的截面.}

事实上, 由结构搬移, \(u(a)\) 此时为 \(\mathfrak g\) 的正则截面, 同时包含在两个 Cartan 子代数 \(\mathfrak d\) 和 \(u(\mathfrak d)\) 中; 所以\SGARefLink{sga3:xiv:2.6}{命题~2.6} 表明这两个子代数相等.

% END inputs/XIV/tex/components/02_expose_XIV_proposition26_corollary27_en.tex

% BEGIN inputs/XIV/tex/components/03_expose_XIV_corollary28_and_nilpotent_rank_en.tex
% Authority: Expo14.pdf, local p. 5 / combined p. 839.
% Continuous English ends immediately before Proposition 2.9.

\paragraph{\textbf{推论 2.8.}}
\SGARefTarget{sga3:xiv:2.8}
\emph{在 \SGARefLink{sga3:xiv:2.4}{2.4} 的假设下, 设 \(\mathfrak d\) 为 \(\mathfrak g\) 的 Cartan 子代数. 对每个满足 \(k(s)\) 无限的 \(s\in S\), 都存在 \(s\) 的一个开邻域 \(V\) 和 \(\mathfrak g\) 在 \(V\) 上的一个正则截面 \(a\), 使得
\[
  \mathfrak d|_V=\Nil\!\left(a,\mathfrak g|_V\right);
\]
等价地, \(a\) 是 \(\mathfrak d|_V\) 的截面.}

事实上, 由于 \(k(s)\) 无限, \(\mathfrak d(s)\) 中包含 \(\mathfrak g(s)\) 的一个正则元素 \(a_0\). 将 \(S\) 换成 \(s\) 的某个开邻域 \(U\) 后, 可以将 \(a_0\) 延拓为 \(\mathfrak d\) 在 \(U\) 上的截面 \(a\). \(\operatorname{ad}(a)\) 在纤维
\[
  \mathfrak g(s)/\mathfrak d(s)
\]
上诱导的自同态为自同构. 将 \(U\) 缩小为一个更小的开邻域 \(V\), 因而可假定 \(\operatorname{ad}(a)_{\mathfrak g/\mathfrak d}\) 为自同构. 从而 \(a\) 为 \(\mathfrak g|_V\) 的一个满足所需条件的拟正则截面.

设 \(\mathfrak g\) 如 \SGARefLink{sga3:xiv:2.4}{2.4}. 考察其 Killing 多项式立即可见, \(S\) 上的函数
\[
  s\longmapsto\mathfrak g(s)\text{ 的幂零秩}
\]
上半连续. 我们特别关心此函数连续, 即局部常值的情形. 下面给出此性质的几个变式.

% END inputs/XIV/tex/components/03_expose_XIV_corollary28_and_nilpotent_rank_en.tex

% BEGIN inputs/XIV/tex/components/04_expose_XIV_proposition29_through_corollary211_en.tex
% Authority: Expo14.pdf, local pp. 5--6 / combined pp. 839--840.
% Continuous English ends immediately before Remark 2.12.

\paragraph{\textbf{命题 2.9.}}
\SGARefTarget{sga3:xiv:2.9}
{\itshape 设 \(S,\mathfrak g\) 如 \SGARefLink{sga3:xiv:2.4}{2.4}. 考虑下列条件:

\smallskip
\noindent
\textup{\((\mathrm C_0)\)}
对 \(s\in S\), \(\mathfrak g(s)\) 的幂零秩是 \(s\) 的局部常值函数.

\smallskip
\noindent
\textup{\((\mathrm C_1)\)}
对 fpqc 拓扑局部地, 存在 \(\mathfrak g\) 的 Cartan 子代数.

\smallskip
\noindent
\textup{\((\mathrm C'_1)\)}
与 \((\mathrm C_1)\) 相同, 将“fpqc 拓扑”换成“平展拓扑”.

\smallskip
\noindent
\textup{\((\mathrm C_2)\)}
条件 \((\mathrm C_0)\) 成立, 且对 \(S\) 上的每个 \(S'\),
\[
  \mathfrak g_{S'}=\mathfrak g\otimes_S S'
\]
的每个拟正则截面均正则.

\smallskip
\noindent
\textup{\((\mathrm C_3)\)}
每个 \(s\in S\) 均有一个开邻域 \(V\), 使得 \(\mathfrak g\) 的 Killing 多项式在其上具有如下形式:
\[
  t^r\bigl(t^{n-r}+c_1t^{n-r-1}+\cdots+c_{n-r}\bigr),
\]
其中对每个 \(s\in V\),
\[
  c_{n-r}(s)\in
  \operatorname{Sym}\!\bigl(\mathfrak g(s)^\vee\bigr)
\]
不为零.

\smallskip
沿用这些记号, 有蕴含关系
\[
  (\mathrm C_3)\Longrightarrow(\mathrm C_2)
  \Longrightarrow(\mathrm C'_1)\Longrightarrow(\mathrm C_1)
  \Longrightarrow(\mathrm C_0),
\]
并且, 当 \(S\) 约化时, 上述五个条件等价.}

还应注意, 这些条件显然在基变换下稳定, 并对 fpqc 拓扑是局部的.

蕴含关系
\[
  (\mathrm C'_1)\Longrightarrow(\mathrm C_1)
  \Longrightarrow(\mathrm C_0)
\]
显然成立, 而当 \(S\) 约化时, 蕴含关系 \((\mathrm C_0)\Longrightarrow(\mathrm C_3)\) 是直接的. 此外记下:

\paragraph{\textbf{推论 2.10.}}
\SGARefTarget{sga3:xiv:2.10}
\emph{假定条件 \((\mathrm C_0)\) 成立. 令 \(U\) 为 \(W(\mathfrak g)\) 中在各自纤维内正则的元素所成的集合. 则 \(U\) 为开集. 特别地, 对 \(\mathfrak g\) 在 \(S\) 上的每个截面 \(a\), 满足 \(a(s)\in\mathfrak g(s)\) 正则的点 \(s\in S\) 所成的集合 \(V\) 为开集.}

事实上, 将第二个断言应用于 \(S\) 上每个基变换 \(S'\) 的 \(\mathfrak g_{S'}\), 即得第一个断言. 对于第二个断言, 可假定 \(S\) 约化; 则 \((\mathrm C_3)\) 成立, 只须考察 \(a\) 在 \(\mathfrak g\) 中的 Killing 多项式即可.

蕴含关系 \((\mathrm C_2)\Longrightarrow(\mathrm C'_1)\) 现可由下述更精确陈述的 b) 部分容易地推出.

\paragraph{\textbf{推论 2.11.}}
\SGARefTarget{sga3:xiv:2.11}
{\itshape 假定条件 \((\mathrm C_2)\) 成立. 则:

\smallskip
\noindent
\textup{a)}
对每个 \(s\in S\) 以及 \(\mathfrak g(s)\) 的每个 Cartan 子代数 \(\mathfrak d_0\), 若 \(\mathfrak d_0\) 包含 \(\mathfrak g(s)\) 的正则元素\textup{（当 \(k(s)\) 无限时, 此条件自动满足）}, 均存在 \(s\) 的一个开邻域 \(V\) 及 \(\mathfrak g|_V\) 的一个 Cartan 子代数 \(\mathfrak d\), 它在 \(s\) 处的纤维为 \(\mathfrak d_0\). 若 \(S_1\) 是 \(S\) 的一个包含 \(s\) 的子概形, 且 \(\mathfrak d_0\) 已延拓为 \(\mathfrak g\otimes_S S_1\) 的 Cartan 子代数 \(\mathfrak d_1\), 则可找到 \(s\) 在 \(S\) 中的一个开邻域 \(V\) 及 \(\mathfrak g|_V\) 的一个 Cartan 子代数 \(\mathfrak d\), 使得
\[
  \mathfrak d\otimes_V(S_1\cap V)
  =\mathfrak d_1|_{S_1\cap V}.
\]

\smallskip
\noindent
\textup{b)}
对每个满足 \(\mathfrak g(s)\) 包含正则元素的 \(s\in S\)\textup{（当 \(k(s)\) 无限时, 此条件自动满足）}, 存在 \(s\) 的一个开邻域 \(V\) 以及 \(\mathfrak g|_V\) 的一个 Cartan 子代数 \(\mathfrak d\).}

在 a) 中取
\[
  \mathfrak d_0=\Nil\!\left(a_0,\mathfrak g(s)\right),
\]
即可得到 b), 其中 \(a_0\) 是 \(\mathfrak g(s)\) 的正则元素.
为证明 a) 的第二种陈述, 选取 \(\mathfrak d_0\) 中的一个在 \(\mathfrak g(s)\) 中正则的元素 \(a_0\). 在 \(S_1\) 中的 \(s\) 附近将它延拓为 \(\mathfrak d_1\) 的截面, 再在 \(s\) 附近将该截面延拓为 \(\mathfrak g\) 的截面. 由 \SGARefLink{sga3:xiv:2.10}{2.10}, 该截面在 \(s\) 的某个开邻域 \(V\) 上拟正则, 因而由 \((\mathrm C_2)\) 它正则. 所以, 由 \SGARefLink{sga3:xiv:2.6}{2.6},
\[
  \mathfrak d=\Nil\!\left(a,\mathfrak g|_V\right)
\]
具有所需性质.

% END inputs/XIV/tex/components/04_expose_XIV_proposition29_through_corollary211_en.tex

% BEGIN inputs/XIV/tex/components/05_expose_XIV_remark212_lemma213_en.tex
% Authority: Expo14.pdf, local p. 7 / combined p. 841.
% Continuous English ends after the proof of Lemma 2.13, hence of 2.9.

\paragraph{\textbf{注 2.12.}}
\SGARefTarget{sga3:xiv:2.12}
我不知道, 当 \(k(s)\) 有限时, 去掉存在正则点的假设后, \SGARefLink{sga3:xiv:2.11}{2.11}~a) 和 b) 的断言是否仍然成立.

尚须证明
\[
  (\mathrm C_3)\Longrightarrow(\mathrm C_2).
\]
还记下 \((\mathrm C_3)\) 的下述等价形式:

\smallskip
\noindent
\textup{\((\mathrm C'_3)\)}
条件 \((\mathrm C_0)\) 成立, 即对 \(s\in S\), \(\mathfrak g(s)\) 的幂零秩局部常值; 并且在 \(S\) 的每个使此秩等于 \(r\) 的开子集 \(V\) 上, \(\mathfrak g|_V\) 的 Killing 多项式可被 \(t^r\) 整除.

\smallskip
须证明此条件蕴含 \(\mathfrak g\) 的每个拟正则截面 \(a\) 均正则. 根据 \SGARefLink{sga3:xiv:2.6}{2.6}, 将下述引理中的蕴含关系 \(\textup{(iv)}\Longrightarrow\textup{(iii)}\) 应用于 \(\mathfrak g\) 的自同态 \(\operatorname{ad}(a)\), 即得此结论.

\paragraph{\textbf{引理 2.13.}}
\SGARefTarget{sga3:xiv:2.13}
{\itshape 设 \(A\) 是环, \(M\) 是有限型投射 \(A\)-模, \(u\) 是 \(M\) 的自同态. 下列条件等价:

\smallskip
\noindent
\textup{(i)}
\(M\) 是两个 \(u\)-稳定子模 \(M'\) 与 \(M''\) 的直和, 满足 \(u|_{M'}\) 幂零, 而 \(u|_{M''}\) 是 \(M''\) 的自同构.

\smallskip
\noindent
\textup{(ii)}
存在整数 \(n>0\), 使得
\[
  \operatorname{Im}u^n+\operatorname{Ker}u^n=M.
\]

\smallskip
\noindent
\textup{(iii)}
幂零空间
\[
  \Nil(u)=\bigcup_{n>0}\operatorname{Ker}u^n
\]
是 \(M\) 的直和因子, 且
\[
  M=\Nil(u)+u(M).
\]

\smallskip
下述条件蕴含这些条件, 并且当 \(A\) 约化时与它们等价:

\smallskip
\noindent
\textup{(iv)}
在 \(\operatorname{Spec}(A)\) 上对 Zariski 拓扑局部地, \(u\) 的特征多项式 \(P_u(t)\) 可写为
\[
  t^r\bigl(t^{n-r}+c_1t^{n-r-1}+\cdots+c_{n-r}\bigr),
\]
其中 \(c_{n-r}\) 可逆.}

\textup{(i)}、\textup{(ii)} 与 \textup{(iii)} 的等价性是直接的, 此处列出以便引用. 当 \(A\) 约化时, \textup{(i)} 蕴含 \textup{(iv)}, 因为此时秩为 \(r\) 的投射模的幂零自同态的特征多项式为 \(t^r\). 另一方面, 在所有情形中, 有限型投射模的自同构的特征多项式的常数项, 在 \(\operatorname{Spec}(A)\) 上局部地, 除去一个符号外就是 \(u\) 的行列式; 因而它是 \(A\) 的一个可逆元.

最后, 为证明 \(\textup{(iv)}\Longrightarrow\textup{(i)}\), 将 \(M\) 看作多项式环 \(A[t]\) 上的模, 其中 \(t\) 以 \(u\) 作用. 熟知的恒等式
\[
  P(u)=0
\]
表明 \(M\) 被 \(P A[t]\) 零化, 因而可看作 \(A[t]/P A[t]\) 上的模. 写成
\[
  P=t^rQ,
\]
其中 \(Q\) 的常数项可逆. 则
\[
  P A[t]=t^rA[t]\cap Q A[t],
\]
所以 \(A[t]/P A[t]\) 分解为下述两个环的乘积:
\[
  A[t]/t^rA[t]
  \qquad\text{与}\qquad
  A[t]/Q A[t].
\]
这相应地给出 \(M\) 的分解, 使之为两个 \(A[t]\)-模的直和, 等价地, 为两个 \(u\)-稳定 \(A\)-子模 \(M'\) 与 \(M''\) 的直和; 这恰好是 \textup{(i)} 中的分解.

这就完成了 \SGARefLink{sga3:xiv:2.13}{2.13} 的证明, 因而也完成了 \SGARefLink{sga3:xiv:2.9}{2.9} 的证明.

% END inputs/XIV/tex/components/05_expose_XIV_remark212_lemma213_en.tex

% BEGIN inputs/XIV/tex/components/06_expose_XIV_corollary214_remark215_en.tex
% Authority: Expo14.pdf, local p. 8 / combined p. 842.
% Continuous English ends immediately after Remark 2.15(c).

\paragraph{\textbf{推论 2.14.}}
\SGARefTarget{sga3:xiv:2.14}
\emph{当 \SGARefLink{sga3:xiv:2.9}{2.9} 的条件 \((\mathrm C_3)\) 成立时, \(\mathfrak g\) 的 Cartan 子代数严格幂零\textup{（\SGARefLink{sga3:xiv:2.2}{2.2}）}.}

证明是直接的.

\paragraph{\textbf{注 2.15.}}
\SGARefTarget{sga3:xiv:2.15}
\noindent
\textup{a)}
可以证明 \SGARefLink{sga3:xiv:2.14}{2.14} 的一个逆命题: 条件 \((\mathrm C_3)\) 等价于如下断言: 对 \(S\) 上的每个 \(S'\), \(\mathfrak g_{S'}\) 的每个拟正则截面均正则, 且 \(\mathfrak g_{S'}\) 的每个 Cartan 子代数均严格幂零; 等价地, \(\mathfrak g_{S'}\) 的每个拟正则截面都包含在 \(\mathfrak g_{S'}\) 的一个严格幂零 Cartan 子代数中.

\smallskip
\noindent
\textup{b)}
与其他条件 \((\mathrm C_0)\) 至 \((\mathrm C_2)\) 不同, 条件 \((\mathrm C_3)\) 具有“无穷小性质”. 更精确地说, 当 \(S\) 局部 Noether 时, \(\mathfrak g\) 满足 \((\mathrm C_3)\) 当且仅当对每个 \(S'\) 为局部 Artin 概形的基变换 \(S'\to S\), \(\mathfrak g_{S'}\) 满足同一条件. 如有需要, 可取 \(S'\) 为 \(S\) 的某个局部环的一个 Artin 商的谱. 同样, 当 \((\mathrm C_0)\) 成立时, \(\mathfrak g\) 的截面为正则的条件也具有无穷小性质.

\smallskip
\noindent
\textup{c)}
当 \(\mathfrak g\) 是 \(S\) 上的有限表示光滑群预概形的 Lie 代数时, 我们将看到 \(\mathfrak g\) 的条件
\[
  (\mathrm C_0),\quad(\mathrm C_1),\quad
  (\mathrm C'_1),\quad(\mathrm C_2)
\]
等价\textup{（\SGARefLink{sga3:xiv:5.2}{5.2}~a）}. 我不知道一般情形如何, 只知道即使 \(S\) 是局部 Artin 概形, \((\mathrm C_0)\) 也不蕴含 \((\mathrm C_1)\). 然而, 即使 \(\mathfrak g\) 来自群 \(G\) 且 \(S\) 是局部 Artin 概形, 一般而言 \((\mathrm C_2)\) 也不蕴含 \((\mathrm C_3)\), 因为 \(\mathfrak g\) 可能幂零而不严格幂零.

一个例子可由离散赋值环的谱 \(S\) 上的光滑仿射群概形 \(G\) 得到, 使一般纤维的 Lie 代数不幂零——例如, 一般纤维可以是伴随型半单群——而特殊纤维的 Lie 代数幂零——例如, 特殊纤维可以是向量群. 此时对充分大的 \(n\),
\[
  G_n=G\times_S S_n
\]
的 Lie 代数虽幂零却不严格幂零.

% END inputs/XIV/tex/components/06_expose_XIV_corollary214_remark215_en.tex

% BEGIN inputs/XIV/tex/components/07_expose_XIV_cartans_functor_theorem216_statement_en.tex
% Authority: Expo14.pdf, local pp. 8--9 / combined pp. 842--843.
% Continuous English ends immediately before the proof of Theorem 2.16.

继续沿用 \SGARefLink{sga3:xiv:2.4}{2.4} 的假设. 令
\[
  \mathcal D\colon
  \bigl((\mathrm{Sch})_{/S}\bigr)^\circ
  \longrightarrow \mathrm{Sets}
\]
为下述定义的函子:
\[
  \mathcal D(S')
  =\left\{
    \mathfrak g_{S'}\text{ 的 Cartan 子代数}
   \right\}.
\]
还引入函子
\[
  X(S')
  =\left\{
    (\mathfrak d,a)\ \middle|\
    \begin{array}{l}
      \mathfrak d\text{ 是 }\mathfrak g_{S'}\text{ 的 Cartan 子代数},\\
      a\text{ 是 }\mathfrak d\text{ 的截面}
    \end{array}
   \right\}.
\]
两个投影
\[
  (\mathfrak d,a)\longmapsto\mathfrak d
  \qquad\text{与}\qquad
  (\mathfrak d,a)\longmapsto a
\]
给出态射
\[
  p\colon X\longrightarrow\mathcal D
  \qquad\text{与}\qquad
  \psi\colon X\longrightarrow W(\mathfrak g).
\]

当 \((\mathrm C_0)\) 成立时, 仍以 \(U\) 表示 \(W(\mathfrak g)\) 的正则点开子集\textup{（\SGARefLink{sga3:xiv:2.10}{2.10}）}. 此时条件 \((\mathrm C_2)\) 等价于说, 由 \(\psi\) 诱导的态射
\[
  \psi^{-1}(U)\longrightarrow U
\]
是同构; 由 \SGARefLink{sga3:xiv:2.6}{2.6}, 首先已知它是单态射.

立即可见,
\[
  p\colon X\longrightarrow\mathcal D
\]
可由向量丛的投影表示. 更明确地说, 对每个对应于 \(\mathfrak g_{S'}\) 的 Cartan 子代数 \(\mathfrak d\) 的 \(S\)-态射 \(S'\to\mathcal D\), 纤维积
\[
  X\times_{\mathcal D}S'
\]
由 \(S'\) 上的向量丛 \(W(\mathfrak d)\) 表示. 因此, 若 \(\mathcal D\) 可表, 则 \(X\) 也可表. 有:

\paragraph{\textbf{定理 2.16.}}
\SGARefTarget{sga3:xiv:2.16}
{\itshape 设 \(S\) 是预概形, \(\mathfrak g\) 是 \(S\) 上的 Lie 代数, 其底层 \(\mathcal O_S\)-模为局部自由有限型模. 假定 \SGARefLink{sga3:xiv:2.9}{2.9} 的条件 \((\mathrm C_0)\) 成立.

\smallskip
\noindent
\textup{a)}
上面定义的 \(\mathfrak g\) 的 Cartan 子代数函子 \(\mathcal D\) 由 \(S\) 上的有限表示拟射影预概形表示. 对函子 \(X\) 亦然.

\smallskip
\noindent
\textup{b)}
若 \SGARefLink{sga3:xiv:2.9}{2.9} 的条件 \((\mathrm C_2)\) 成立, 则 \(\mathcal D\) 与 \(X\) 在 \(S\) 上光滑, 且由 \(\psi\) 诱导的态射
\[
  \psi^{-1}(U)\longrightarrow U
\]
是同构.

\smallskip
\noindent
\textup{c)}
继续假定 \((\mathrm C_2)\) 成立. 设 \(s\in S\), \(\mathfrak d_0\) 是 \(\mathfrak g(s)\) 的一个 Cartan 子代数, 对应于 \(\mathcal D\) 的一个 \(k(s)\)-有理点 \(d\). 假定 \(\mathfrak d_0\) 包含 \(\mathfrak g(s)\) 的正则元素, 当 \(k(s)\) 无限时此条件自动满足. 若 \(r\) 是 \(\mathfrak g(s)\) 的无穷小秩, \(n\) 是其在 \(k(s)\) 上的秩, 则 \(d\) 在 \(\mathcal D\) 中存在一个开邻域 \(V\), 它与下式的一个开子集 \(V'\) \(S\)-同构:
\[
  S[t_1,\ldots,t_{n-r}].
\]}

% END inputs/XIV/tex/components/07_expose_XIV_cartans_functor_theorem216_statement_en.tex

% BEGIN inputs/XIV/tex/components/08_expose_XIV_theorem216_proof_corollary217_remark218_en.tex
% Authority: Expo14.pdf, local pp. 9--11 / combined pp. 843--845.
% Continuous English ends at the close of Section 2, before Section 3.

\paragraph{\textit{证明.}}
可假定 \(\mathfrak g\) 的秩为常数 \(n\), 幂零秩为常数 \(r\). \textup{a)} 与 \textup{b)} 中关于 \(X\) 的断言, 由关于 \(\mathcal D\) 的断言以及 \(X\) 为 \(\mathcal D\) 上由一个局部自由模定义的向量丛这一事实立即得到; 将它们列出仅为使陈述完整.

\medskip
\noindent\textup{a)}
函子 \(\mathcal D\) 是 \(\operatorname{Grass}_{n-r}(\mathfrak g)\) 的子函子, 后者在 \(S'\) 处的值是 \(\mathfrak g_{S'}\) 的秩为 \(n-r\) 的局部自由商模所成的集合. 众所周知, 后一函子由 \(S\) 上光滑的射影预概形表示; 例如参见 Cartan Seminar 1960/61, 第~12 讲, 第~2、3 节, 其中的构造可原样推广到预概形.\begingroup
\renewcommand{\thefootnote}{*}\footnote{亦参见 EGA~I, 第二版\textup{（当时宣布将由 North-Holland Publishing Company 出版）}.}
\endgroup\footnote{编者注: 参见 M.~Demazure 与 P.~Gabriel, \emph{Groupes algébriques}, Masson \textup{(1970)}, \S \SGARefLink{sga3:xiv:1.3}{I.1.3}.}

因此归结到下述相对问题. 给定 \(\mathfrak g\) 的秩为 \(n-r\) 的局部自由商模, 等价地, 给定一个秩为 \(r\) 且局部为直和因子的局部自由子模 \(\mathfrak d\), 表示函子
\[
 F(S')=
 \begin{cases}
   \varnothing,
     &\text{若 \(\mathfrak d_{S'}\) 不是 \(\mathfrak g_{S'}\) 的 Cartan 子代数,}\\
   \{\varnothing\},
     &\text{其他情形.}
 \end{cases}
\]
将证明 \(F\) 由 \(S\) 上有限表示的子预概形表示. 这就表明
\[
  \mathcal D\longrightarrow
  \operatorname{Grass}_{n-r}(\mathfrak g)
\]
由有限表示的浸入表示, 从而完成 a) 的证明.

首先表达 \(\mathfrak d_{S'}\) 为 \(\mathfrak g_{S'}\) 的 Lie 子代数这一条件. 此条件恰好就是说 \(S'\to S\) 通过 \(S\) 上的某个有限表示闭子预概形 \(S_1\) 分解; 在 \(S\) 上局部地, 从 \(\mathfrak g\) 的一个适合于 \(\mathfrak d\) 的基立即可得其方程. 因而可假定 \(S=S_1\).

其次须表达 \(\mathfrak d_{S'}\) 对 fpqc 拓扑局部地包含 \(\mathfrak g_{S'}\) 的拟正则截面这一条件. 为此考虑
\[
  V=W(\mathfrak d)\cap U,
\]
其中 \(U\) 是 \(W(\mathfrak g)\) 的正则点开子集\textup{（\SGARefLink{sga3:xiv:2.10}{2.10}）}. 因为结构态射 \(V\to S\) 光滑且拟紧, 其像 \(S_1\) 是 \(S\) 的开子集, 且浸入 \(S_1\to S\) 拟紧, 因而为有限表示. 对 \(S'\) 的条件恰好是 \(S'\to S\) 通过 \(S_1\) 分解. 于是归结到 \(S=S_1\) 的情形; 再由下降归结到 \(\mathfrak d\) 有一个在 \(\mathfrak g\) 中拟正则的截面 \(a\) 的情形.

只须再表达, 对所诱导的截面 \(a_{S'}\), 自同态
\[
  \operatorname{ad}(a_{S'})
    _{\mathfrak g_{S'}/\mathfrak d_{S'}}
\]
为双射. 这仍然就是说 \(S'\to S\) 通过某个有限表示开子预概形分解, 即通过 \(S_D\) 分解, 其中
\[
  D=\det\!\left(
    \operatorname{ad}(a)_{\mathfrak g/\mathfrak d}
  \right).
\]
而此时 \(\mathfrak d|_{S_D}\) 为 \(\mathfrak g|_{S_D}\) 的 Cartan 子代数. 因此 \(S_D\) 表示 \(F\), 这就证明了 a).

\medskip
\noindent\textup{b)}
这由 \SGARefLink{sga3:xiv:2.11}{2.11}\textup{(a)} 和 \SGARefLink{sga3:xi:remark:1:5}{XI~1.5} 立即得到. 当然, 它也是更精确的断言 c) 的推论.

\medskip
\noindent\textup{c)}
令 \(a_0\) 为包含在 \(\mathfrak d_0\) 中的 \(\mathfrak g(s)\) 的正则点, 并将其延拓为 \(\mathfrak g\) 在 \(s\) 的某个开邻域上的截面 \(a\). 显然可假定此邻域为整个 \(S\). 在向量空间 \(\mathfrak g(s)\) 中选取 \(\mathfrak d_0\) 的一个补空间 \(M_0\). 再将 \(S\) 换成 \(s\) 的一个开邻域, 可找到 \(\mathfrak g\) 的一个局部为直和因子的子模 \(M\), 使得 \(M(s)=M_0\); 仍假定此邻域为 \(S\).

令 \(\mathcal V\) 为 \(\mathcal D\) 的子函子, 其中 \(\mathcal V(S')\) 由满足下述条件的 \(\mathfrak g_{S'}\) 的 Cartan 子代数 \(\mathfrak d'\) 组成:
\begin{enumerate}
\item[\textup{1\textdegree)}]
  \(\mathfrak d'\) 是 \(M_{S'}\) 的补模;
\item[\textup{2\textdegree)}]
  \((a_{S'}+M_{S'})\cap\mathfrak d'\) 中的唯一截面是 \(\mathfrak g_{S'}\) 的正则截面.
\end{enumerate}
条件 \textup{1\textdegree)} 定义 \(\mathcal D\) 的开集 \(V_1\), 它由 \(\operatorname{Grass}_{n-r}(\mathfrak g)\) 中同一条件给出的开集所诱导. 由 \SGARefLink{sga3:xiv:2.10}{2.10} 和 \((\mathrm C_2)\), \textup{1\textdegree)} 与 \textup{2\textdegree)} 的合取定义 \(V_1\) 的一个开集. 因此 \(\mathcal V\) 由 \(\mathcal D\) 的一个开子预概形 \(V\) 表示, 显然它包含点 \(d\).

再令 \(\mathcal V'\) 为 \(W(M)\) 的如下子函子:
\[
 \mathcal V'(S')
 =
 \left\{
   u'\in\Gamma(S',M_{S'})
   \ \middle|\
   \begin{array}{l}
   \textup{(i) }a_{S'}+u'\text{ 是 }\mathfrak g_{S'}\text{ 的正则截面},\\
   \textup{(ii) 唯一的 Cartan 子代数 }\mathfrak d'
     \text{ 属于 }\mathfrak g_{S'}\\
   \quad\text{并包含 }a_{S'}+u'
     \text{, 是 }M_{S'}\text{ 的补模}
   \end{array}
 \right\}.
\]
条件 \textup{(i)} 定义 \(W(M)\) 的一个开子预概形 \(V'_1\): 它是 \(W(\mathfrak g)\) 的正则点开集 \(U\)\textup{（\SGARefLink{sga3:xiv:2.10}{2.10}）} 在平移态射 \(m\mapsto a+m\) 下的逆像. \textup{(i)} 与 \textup{(ii)} 的合取定义 \(V'_1\) 的一个开集 \(V'\), 即 \(V\) 在显然的态射
\[
  V'_1\longrightarrow\mathcal D
\]
下的逆像, 该态射将 \(u'\) 对应到唯一的包含 \(a_{S'}+u'\) 的 Cartan 子代数 \(\mathfrak d'\). 它的限制给出态射
\[
  V'\longrightarrow V,
\]
它显然是同构. 这就证明了 c).

\paragraph{\textbf{推论 2.17.}}
\SGARefTarget{sga3:xiv:2.17}
{\itshape 设 \(\mathfrak g\) 是域 \(k\) 上的有限维 Lie 代数. 则 \(\mathfrak g\) 的 Cartan 子代数概形 \(\mathcal D\)\textup{（\SGARefLink{sga3:xiv:2.16}{2.16}\textup{(a)}）} 拟射影、光滑且不可约. 若 \(\mathfrak g\) 包含一个正则元素\textup{（例如, 若 \(k\) 无限）}, 则 \(\mathcal D\) 为有理簇; 即其函数域为 \(k\) 的纯超越扩张.}

事实上, \(\mathcal D\) 不可约, 因为存在满态射
\[
  \psi^{-1}(U)\longrightarrow\mathcal D,
\]
且 \(\psi^{-1}(U)\) 不可约, 因为它同构于 \(W(\mathfrak g)\) 的开子集 \(U\). 关于函数域的断言由 \SGARefLink{sga3:xiv:2.16}{2.16}\textup{(c)} 立即得到.

\paragraph{\textbf{注 2.18.}}
\SGARefTarget{sga3:xiv:2.18}
我不知道, 当 \(k\) 有限时, 不假定 \(\mathfrak g\) 包含正则点, 这一结论是否仍成立; 比较 \SGARefLink{sga3:xiv:2.12}{2.12}. 当 \(\mathfrak g\) 是 \(k\) 上光滑代数群 \(G\) 的 Lie 代数时, 至少当 \(G/\operatorname{radical}(G)\) 为“伴随型”半单群时, 可以利用下文所引 Chevalley 的一个结果\textup{（参见附录）} 证明它. 看来即使不对 \(G\) 加以限制, 结果仍应成立; 为此只须对每个未必为伴随型的半单代数群证明所引 Chevalley 的结果.

% END inputs/XIV/tex/components/08_expose_XIV_theorem216_proof_corollary217_remark218_en.tex

% BEGIN inputs/XIV/tex/components/09_expose_XIV_theorem31_en.tex
% Authority: Expo14.pdf, local pp. 11--12 / combined pp. 845--846.
% Continuous English ends after the proof of Theorem 3.1.

\section{任意预概形上的群预概形的 \textup{(C)} 型子群}
\SGARefTarget{sga3:xiv:section:3}

\paragraph{\textbf{定理 3.1.}}
\SGARefTarget{sga3:xiv:3.1}
{\itshape 设 \(S\) 是预概形, \(G\) 是光滑 \(S\)-群预概形, \(\mathfrak g\) 是它的 Lie 代数, 它是 \(S\) 上的局部自由有限型模. 设 \(\mathfrak h\) 是 \(\mathfrak g\) 的 Lie 子代数, 作为模在 \(\mathfrak g\) 中局部为直和因子; 假定对每个 \(s\in S\), 几何纤维 \(\mathfrak h_{\bar s}\) 包含 \(\mathfrak g_{\bar s}\) 的一个 Cartan 子代数. 设 \(a\) 是 \(\mathfrak g\) 的拟正则截面\textup{（\SGARefLink{sga3:xiv:2.5}{2.5}）}. 则
\[
  M_a=\operatorname{Transp}_G(a,\mathfrak h),
\]
即 \(G\) 的一个子函子, 它的 \(S'\)-值点为满足下式的 \(g\in G(S')\):
\[
  \operatorname{ad}(g)\cdot a_{S'}
    \in\Gamma(S',\mathfrak h_{S'}),
\]
由 \(G\) 的一个在 \(S\) 上光滑的闭子预概形表示, 且其到 \(S\) 的结构态射为满射.}
\footnote{编者注: 此搬运子记为 \(\operatorname{Transp}_G(a,\mathfrak h)\).}

\paragraph{\textit{证明.}}
考虑典范态射
\[
  \varphi\colon G\times_S W(\mathfrak h)\longrightarrow
  W(\mathfrak g),
  \qquad
  (g,x)\longmapsto\operatorname{ad}(g)\cdot x.
\]
则 \(M_a\) 与 \(\varphi^{-1}(a)\) \(S\)-同构, 后者是将 \(a\) 看作 \(W(\mathfrak g)\) 在 \(S\) 上的截面时, 它在 \(\varphi\) 下的逆像. 为证明 \(M_a\) 光滑, 只须表明 \(\varphi\) 在 \(G\times_S W(\mathfrak h)\) 中位于 \(\operatorname{Im}(a)\) 上方的各点处光滑. 更一般地, \(\varphi\) 在所有位于 \(W(\mathfrak g)\) 中的某个在其于 \(S\) 上的纤维 \(W(\mathfrak g(s))\) 中正则的点上方的点处光滑.

事实上, \(\varphi\) 的源和目标均在 \(S\) 上光滑, 因而在 \(S\) 上平坦且局部为有限表示. 所以可逐纤维检验此断言. 这将我们归结到如下情形: \(S\) 是代数闭域 \(k\) 的谱, \(G\) 是 \(k\) 上的局部代数群, \(\mathfrak h\) 是其 Lie 代数 \(\mathfrak g\) 的一个包含 \(\mathfrak g\) 的 Cartan 子代数的 Lie 子代数, 而 \(a\) 是 \(\mathfrak g\) 的正则点.

由于 \(\varphi\) 为 \(G\)-态射, 可假定所考虑的 \(G\times W(\mathfrak h)\) 的点形如 \((e,a)\). 也可假定 \(G\) 连通, 因而在 \(k\) 上为有限型. 此时断言恰好就是 \SGARefLink{sga3:xiii:corollary:5:4}{XIII~5.4}.

\(M_a\) 为 \(G\) 的一个在 \(S\) 上有限表示的闭子预概形也是直接的: 它是 \(W(\mathfrak h)\) 在态射
\[
  G\longrightarrow W(\mathfrak g),
  \qquad
  g\longmapsto\operatorname{ad}(g)\cdot a.
\]
下的逆像. 最后, \(M_a\to S\) 的满射性同样归结到基域代数闭的情形. 根据假设, 此时 \(\mathfrak h\) 包含一个 Cartan 子代数 \(\mathfrak d'\). 根据共轭定理 \SGARefLink{sga3:xiii:theorem:6:1}{XIII~6.1}\textup{(a)}, \(\mathfrak d'\) 与
\[
  \mathfrak d=\Nil(a,\mathfrak g).
\]
共轭. 从而 \(\mathfrak h\) 包含 \(a\) 的某个共轭元素. 这就完成证明.

% END inputs/XIV/tex/components/09_expose_XIV_theorem31_en.tex

% BEGIN inputs/XIV/tex/components/10_expose_XIV_corollary32_lemma33_en.tex
% Authority: Expo14.pdf, local p. 12 / combined p. 846.
% Continuous English ends immediately before Corollary 3.4.

\paragraph{\textbf{推论 3.2.}}
\SGARefTarget{sga3:xiv:3.2}
{\itshape 设 \(G\) 与 \(\mathfrak g\) 如 \SGARefLink{sga3:xiv:3.1}{3.1}, 其中 \(G\) 在 \(S\) 上为有限型; 设 \(\mathfrak k\) 与 \(\mathfrak h\) 为 \(\mathfrak g\) 的两个 Lie 子代数, 作为模局部为直和因子. 假定下列两个假设之一成立:

\smallskip
\noindent
\textup{a)}
对每个 \(s\in S\), 几何纤维 \(\mathfrak k_{\bar s}\) 幂零且包含 \(\mathfrak g_{\bar s}\) 的一个正则元素, 而几何纤维 \(\mathfrak h_{\bar s}\) 包含 \(\mathfrak g_{\bar s}\) 的一个 Cartan 子代数.

\smallskip
\noindent
\textup{b)}
\(\mathfrak k\) 是 \(\mathfrak g\) 的 Cartan 子代数.

\smallskip
\noindent
则 \(\operatorname{Transp}_G(\mathfrak k,\mathfrak h)\) 是 \(G\) 的一个在 \(S\) 上光滑的闭子预概形. 此外, 在情形 a) 中, 它到 \(S\) 的结构态射为满射.}

搬运子可由 \(G\) 的一个有限表示闭子预概形表示是直接的, 留给读者. 先证明情形 a) 中的光滑性. 首先假定 \(\mathfrak k\) 有一个在 \(\mathfrak g\) 中拟正则的截面 \(a\). 此时只须应用 \SGARefLink{sga3:xiv:3.1}{3.1} 和下述引理.

\paragraph{\textbf{引理 3.3.}}
\SGARefTarget{sga3:xiv:3.3}
{\itshape 在 \SGARefLink{sga3:xiv:3.2}{3.2}\textup{(a)} 的假设下, 若 \(a\) 是 \(\mathfrak k\) 的一个在 \(\mathfrak g\) 中拟正则的截面, 则
\[
  \operatorname{Transp}_G(\mathfrak k,\mathfrak h)
  =
  \operatorname{Transp}_G(a,\mathfrak h).
\]}

事实上, 由定义只须表明, 若 \(a\) 也是 \(\mathfrak h\) 的截面, 则 \(\mathfrak k\subset\mathfrak h\). 根据假设, \(\mathfrak k\) 局部幂零, 故 \SGARefLink{sga3:xiv:2.1}{2.1} 给出
\[
  \mathfrak k\subset\Nil(a,\mathfrak g).
\]
另一方面,
\[
  \Nil(a,\mathfrak g)\subset\mathfrak h,
\]
因为
\[
  \operatorname{ad}(a)_{\mathfrak g/\mathfrak h}
\]
为单射; 根据 \SGARefLink{sga3:xiii:proposition:4:8:item:b}{XIII~4.8(b)}, 可逐纤维检验这一点. 引理得证.

一般地, 通常的标准论证先将问题归结到 \(S\) 为仿射 Noether 概形的情形, 然后利用光滑性是无穷小性质, 进一步归结到 \(S\) 为局部 Artin 概形的情形. 由平坦下降, 还可假定剩余域无限. 此时 \(\mathfrak k\) 的纤维包含一个在 \(\mathfrak g\) 的相应纤维中正则的元素. 将此元素提升为 \(\mathfrak k\) 的截面, 就归结到前一种情形.

所以在情形 a) 中搬运子光滑. 其结构态射为满射, 因为这一点可以在归结到 \(S\) 是代数闭域的谱后检验; 此时 \(\mathfrak k\) 包含 \(\mathfrak g\) 的一个正则点, 可应用 \SGARefLink{sga3:xiv:3.3}{3.3} 和 \SGARefLink{sga3:xiv:3.1}{3.1}.

对于情形 b), 根据光滑性的定义\textup{（\SGARefLink{sga3:xi:definition:1:1}{XI~1.1}）}, 断言归结到下述提升命题. 设 \(S\) 仿射, \(S_0\) 是由拟凝聚幂零理想 \(J\) 定义的闭子概形, \(g_0\in G(S_0)\) 将 \(\mathfrak k_0\) 搬入 \(\mathfrak h_0\). 则 \(g_0\) 可提升为某个将 \(\mathfrak k\) 搬入 \(\mathfrak h\) 的元素 \(g\in G(S)\). 关于 \(g_0\) 的假设使这一提升问题落入已证的情形 a). 证明完成.

当然, 若在 \SGARefLink{sga3:xiv:3.2}{3.2}\textup{(b)} 中, 子代数 \(\mathfrak h\) 满足 \SGARefLink{sga3:xiv:3.1}{3.1} 中更强的假设, 则且仅则结构态射
\[
  \operatorname{Transp}_G(\mathfrak k,\mathfrak h)
  \longrightarrow S
\]
为满射. 利用 Hensel 引理 \SGARefLink{sga3:xi:proposition:1:10}{XI~1.10}, 从 \SGARefLink{sga3:xiv:3.1}{3.1} 和 \SGARefLink{sga3:xiv:3.2}{3.2} 得到下述推论.

% END inputs/XIV/tex/components/10_expose_XIV_corollary32_lemma33_en.tex

% BEGIN inputs/XIV/tex/components/11_expose_XIV_corollary34_through_definition38_en.tex
% Authority: Expo14.pdf, local pp. 12--13 / combined pp. 846--847.
% Continuous English ends immediately before Theorem 3.9.

\paragraph{\textbf{推论 3.4.}}
\SGARefTarget{sga3:xiv:3.4}
{\itshape 对 \(G\) 与 \(\mathfrak h\) 作 \SGARefLink{sga3:xiv:3.1}{3.1} 中的假设, 并假定 \(G\) 在 \(S\) 上为有限型, 则下述断言成立:

\smallskip
\noindent
\textup{a)}
对 \(\mathfrak g\) 的每个拟正则截面 \(a\), 平展局部地存在 \(a\) 的一个共轭截面, 它是 \(\mathfrak h\) 的截面.

\smallskip
\noindent
\textup{b)}
对 \(\mathfrak g\) 的每个 Cartan 子代数 \(\mathfrak d\), 平展局部地, \(\mathfrak d\) 与 \(\mathfrak h\) 的一个子代数共轭.}

特别地, 当 \(\mathfrak h\) 本身是 \(\mathfrak g\) 的 Cartan 子代数时, 得到下述结果.

\paragraph{\textbf{推论 3.5.}}
\SGARefTarget{sga3:xiv:3.5}
{\itshape 设 \(G\) 为光滑有限型 \(S\)-群预概形, \(\mathfrak g\) 为其 Lie 代数, \(\mathfrak d\) 与 \(\mathfrak d'\) 为 \(\mathfrak g\) 的两个 Cartan 子代数. 则 \(\operatorname{Transp}_G(\mathfrak d,\mathfrak d')\) 就是将 \(\mathfrak d\) 搬入 \(\mathfrak d'\) 的严格搬运子, 并且为 \(G\) 的一个在 \(S\) 上光滑的闭子预概形, 其结构态射为满射. 平展局部地, \(\mathfrak d\) 与 \(\mathfrak d'\) 共轭.}

此处搬运子与严格搬运子相同, 仅仅因为 \(\mathfrak d\) 与 \(\mathfrak d'\) 在 \(\mathfrak g\) 中局部为直和因子, 且在每点秩相同. 所以 \SGARefLink{sga3:xiv:3.5}{3.5} 是 \SGARefLink{sga3:xiv:3.4}{3.4} 的特例. 取 \(\mathfrak d=\mathfrak d'\) 即得下述结果.

\paragraph{\textbf{推论 3.6.}}
\SGARefTarget{sga3:xiv:3.6}
{\itshape 设 \(G\) 如 \SGARefLink{sga3:xiv:3.5}{3.5}, \(\mathfrak d\) 是 \(\mathfrak g\) 的 Cartan 子代数. 则 \(\operatorname{Norm}_G(\mathfrak d)\) 是 \(G\) 的一个在 \(S\) 上光滑的闭子群预概形, 其 Lie 代数为 \(\mathfrak d\).}

事实上, 最后一个断言就是说 \(\mathfrak d\) 在 \(\mathfrak g\) 中等于其自身的正规化子, 而这由逐纤维成立立即得到.

\paragraph{\textbf{推论 3.7.}}
\SGARefTarget{sga3:xiv:3.7}
{\itshape 设 \(G\) 与 \(\mathfrak g\) 如 \SGARefLink{sga3:xiv:3.5}{3.5}. 则 \SGARefLink{sga3:xiv:2.9}{2.9} 的条件 \((\mathrm C_2)\)、\((\mathrm C'_1)\) 与 \((\mathrm C_1)\) 等价. 换言之, 若 \(\mathfrak g\) 对 fpqc 拓扑局部容许 Cartan 子代数, 则 \(\mathfrak g\) 的每个拟正则截面均正则.}

事实上, 设 \(a\) 为拟正则截面; 证明它正则. 由于问题对 fpqc 拓扑是局部的, 可假定 \(\mathfrak g\) 容许 Cartan 子代数 \(\mathfrak d\). 由 \SGARefLink{sga3:xiv:3.4}{3.4}\textup{(a)}, 平展局部地, \(a\) 与 \(\mathfrak d\) 的一个截面共轭. 于是归结到 \(a\) 是 \(\mathfrak d\) 的截面的情形, 此时由定义立即得到结论.

\paragraph{\textbf{定义 3.8.}}
\SGARefTarget{sga3:xiv:3.8}
{\itshape 设 \(G\) 为预概形 \(S\) 上的光滑群预概形. \(G\) 的一个 \emph{\((\mathrm C)\) 型子群}是一个在 \(S\) 上光滑且具有连通纤维的子群预概形 \(D\subset G\), 满足
\[
  \mathfrak d=\operatorname{Lie}(D)
\]
是 \(\mathfrak g=\operatorname{Lie}(G)\) 的 Cartan 子代数\textup{（\SGARefLink{sga3:xiv:2.4}{2.4}）}. 等价地, 对每个 \(s\in S\), 纤维 \(D_s\) 是代数群 \(G_s\) 的一个 \((\mathrm C)\) 型子群\textup{（\SGARefLink{sga3:xiii:definition:6:2}{XIII~6.2}）}.}

% END inputs/XIV/tex/components/11_expose_XIV_corollary34_through_definition38_en.tex

% BEGIN inputs/XIV/tex/components/12_expose_XIV_theorem39_through_corollary311_en.tex
% Authority: Expo14.pdf, local pp. 13--14 / combined pp. 847--848.
% Continuous English ends immediately before Proposition 3.12.

\paragraph{\textbf{定理 3.9.}}
\SGARefTarget{sga3:xiv:3.9}
{\itshape 设 \(G\) 为有限表示光滑 \(S\)-群预概形, \(\mathfrak g\) 为其 Lie 代数. 则:

\smallskip
\noindent
\textup{a)}
映射
\[
  D\longmapsto\mathfrak d=\operatorname{Lie}(D)
\]
在 \(G\) 的 \((\mathrm C)\) 型子群与 \(\mathfrak g\) 的 Cartan 子代数之间建立一一对应.

\smallskip
\noindent
\textup{b)}
若 \(D\) 与 \(\mathfrak d\) 相对应, 则
\[
  \operatorname{Norm}_G(D)=\operatorname{Norm}_G(\mathfrak d).
\]
这是 \(G\) 的一个在 \(S\) 上光滑的闭子预概形, 并且
\[
  D=\operatorname{Norm}_G(D)^0
   =\operatorname{Norm}_G(\mathfrak d)^0.
\]

\smallskip
\noindent
\textup{c)}
\(G\) 中任意两个 \((\mathrm C)\) 型子群 \(D\) 与 \(D'\) 平展局部共轭.}

\emph{证明.}
设 \(D\) 为 \(G\) 的一个 \((\mathrm C)\) 型子群, \(\mathfrak d\) 为其 Lie 代数. 则
\[
  D\subset\operatorname{Norm}_G(\mathfrak d).
\]
由\SGARefLink{sga3:xiv:3.8}{定义 3.8} 与\SGARefLink{sga3:xiv:3.6}{推论 3.6}, 这是 \(S\) 上的光滑群预概形之间的包含关系, 并在 Lie 代数上诱导同构. 由于 \(D\) 具有连通纤维, 得到
\[
  D=\operatorname{Norm}_G(\mathfrak d)^0.
\]
因此 a) 中的映射为单射. 下面证明它为满射. 设 \(\mathfrak d\) 为 \(\mathfrak g\) 的 Cartan 子代数. 由 \SGARefLink{sga3:xiv:3.6}{3.6},
\[
  N=\operatorname{Norm}_G(\mathfrak d)
\]
是 \(G\) 的一个在 \(S\) 上光滑的闭子群预概形, 且其 Lie 代数为 \(\mathfrak d\). 由于 \(G\) 在 \(S\) 上为有限表示, \(N\) 亦然. 所以, 如 XII 中 \SGARefLink{sga3:xii:7.3}{7.3} 之后所指出的, \(N\) 的各纤维的连通分支的并是一个开子群预概形 \(N^\circ\subset N\) 的底层集合. 显然 \(N^\circ\) 是 \(G\) 的一个 Lie 代数为 \(\mathfrak d\) 的 \((\mathrm C)\) 型子群. 这就证明了 a). b) 的第一个断言立即得到, 而公式
\[
  D=\operatorname{Norm}_G(\mathfrak d)^0
\]
已经证明. 最后, c) 来自 a) 与 \SGARefLink{sga3:xiv:3.5}{3.5}.

\paragraph{\textbf{推论 3.10.}}
\SGARefTarget{sga3:xiv:3.10}
{\itshape 假定 \(\mathfrak g\) 对 fpqc 拓扑局部容许 Cartan 子代数, 等价地, 根据 \SGARefLink{sga3:xiv:3.9}{3.9}\textup{(a)}, \(G\) 对 fpqc 拓扑局部容许 \((\mathrm C)\) 型子群. 考虑函子
\[
  \mathcal D\colon(\mathrm{Sch})^\circ/S\longrightarrow
  (\mathrm{Sets}),
\]
其中 \(\mathcal D(S')\) 为 \(G_{S'}\) 的 \((\mathrm C)\) 型子群所成的集合. 此函子可由一个在 \(S\) 上光滑且几何纤维连通的拟射影预概形表示. 若 \(S=\operatorname{Spec}(k)\), 从而 \(G\) 为 \(k\) 上的光滑代数群, 并且若 \(\mathfrak g\) 容许正则点——当 \(k\) 无限时此条件自动满足——则 \(\mathcal D\) 是 \(k\) 上的有理簇.}

事实上, 根据 \SGARefLink{sga3:xiv:3.9}{3.9}\textup{(a)}, 函子 \(\mathcal D\) 典范同构于 \SGARefLink{sga3:xiv:2.16}{2.16} 中考虑的函子. 另一方面, 根据 \SGARefLink{sga3:xiv:3.7}{3.7}, 条件 \((\mathrm C_2)\) 满足. 因此 \SGARefLink{sga3:xiv:3.10}{3.10} 来自 \SGARefLink{sga3:xiv:2.16}{2.16} 和 \SGARefLink{sga3:xiv:2.17}{2.17}.

\paragraph{\textbf{推论 3.11.}}
\SGARefTarget{sga3:xiv:3.11}
{\itshape 设 \(D\) 为 \(G\) 的 \((\mathrm C)\) 型子群, \(N\) 为它在 \(G\) 中的正规化子. 则商层 \(G/N\) 与 \SGARefLink{sga3:xiv:3.10}{3.10} 的函子 \(\mathcal D\) 典范同构. 因而它可由一个在 \(S\) 上光滑且几何纤维连通的拟射影预概形表示.}

% END inputs/XIV/tex/components/12_expose_XIV_theorem39_through_corollary311_en.tex

% BEGIN inputs/XIV/tex/components/13_expose_XIV_proposition312_through_corollary315_en.tex
% Authority: Expo14.pdf, local pp. 14--15 / combined pp. 848--849.
% Continuous English ends immediately before Corollary 3.16.

\paragraph{\textbf{命题 3.12.}}
\SGARefTarget{sga3:xiv:3.12}
{\itshape 设 \(G\) 为 \(S\) 上有限表示的光滑 \(S\)-群预概形. 设 \(H\) 与 \(K\) 为两个在 \(S\) 上光滑且有限表示的子群预概形, 其中 \(K\) 具有连通纤维; 以 \(\mathfrak g,\mathfrak h,\mathfrak k\) 表示相应的 Lie 代数. 假定它们的几何纤维满足下列条件之一:

\smallskip
\noindent
\textup{a)}
对每个 \(s\in S\), \(\mathfrak k_{\bar s}\) 包含 \(\mathfrak g_{\bar s}\) 的一个正则元素, 而 \(\mathfrak h_{\bar s}\) 包含 \(\mathfrak g_{\bar s}\) 的一个 Cartan 子代数.

\smallskip
\noindent
\textup{b)}
对每个 \(s\in S\), \(\mathfrak k_{\bar s}\) 包含 \(\mathfrak g_{\bar s}\) 的一个 Cartan 子代数.

\smallskip
\noindent
在这些条件下, \(H\supset K\) 当且仅当 \(\mathfrak h\supset\mathfrak k\).}

只须证明 \(\mathfrak h\supset\mathfrak k\) 蕴含 \(H\supset K\). 在情形 b) 中, 包含关系 \(\mathfrak h\supset\mathfrak k\) 使我们处于情形 a) 的假设下, 所以只须证明 a). 仿照 \SGARefLink{sga3:xiv:3.2}{3.2}\textup{(a)}, 通过归结到 \(S\) 为局部 Artin 概形的情形, 我们归结到存在 \(\mathfrak k\) 的一个在 \(\mathfrak g\) 中拟正则的截面 \(a\) 的情形. 再仿照 \SGARefLink{sga3:xiii:corollary:5:5}{XIII~5.5} 论证, 就归结到下述命题.

\paragraph{\textbf{推论 3.13.}}
\SGARefTarget{sga3:xiv:3.13}
{\itshape 设 \(G\) 为 \(S\) 上有限表示的光滑 \(S\)-群预概形, \(H\subset G\) 是在 \(S\) 上光滑且有限表示的子群预概形; 设 \(\mathfrak g\) 与 \(\mathfrak h\) 为它们的 Lie 代数. 设 \(a\) 是 \(\mathfrak h\) 的一个在 \(\mathfrak g\) 中拟正则的截面. 假定对每个 \(s\in S\), 几何纤维 \(\mathfrak h_{\bar s}\) 包含 \(\mathfrak g_{\bar s}\) 的一个 Cartan 子代数. 置
\[
  M_a=\operatorname{Transp}_G(a,\mathfrak h).
\]
这是 \(G\) 的一个在 \(S\) 上光滑的闭子预概形\textup{（参见 \SGARefLink{sga3:xiv:3.1}{3.1}）}. 因而, \(M_a^0\), 即 \(M_a\) 各纤维中包含单位元的连通分支之并, 是 \(M_a\) 的一个开子集, 赋予由 \(M_a\) 诱导的结构. 则
\[
  H^0=M_a^0.
\]}

显然 \(H\subset M_a\), 从而 \(H^0\subset M_a^0\). 由于这是 \(S\) 上的光滑预概形之间的包含关系, 只须在各纤维上检验相等性. 这归结到 \(S\) 是代数闭域的谱的情形, 该情形已在 \SGARefLink{sga3:xiii:corollary:5:4}{XIII~5.4} 中处理.

\paragraph{\textbf{推论 3.14.}}
\SGARefTarget{sga3:xiv:3.14}
{\itshape 设 \(G\) 为有限表示光滑 \(S\)-群预概形, \(D\) 为 \(G\) 的一个 \((\mathrm C)\) 型子群, \(H\subset G\) 为在 \(S\) 上光滑且有限表示的子群预概形. 若 \(\mathfrak d\) 与 \(\mathfrak h\) 为它们的 Lie 代数, 则
\[
  \operatorname{Transp}_G(D,H)
  =
  \operatorname{Transp}_G(\mathfrak d,\mathfrak h),
\]
且此函子可由 \(G\) 的一个在 \(S\) 上光滑的闭子预概形表示.}

事实上, 两个搬运子相同来自 \SGARefLink{sga3:xiv:3.12}{3.12}\textup{(b)}, 然后可应用 \SGARefLink{sga3:xiv:3.2}{3.2}.

\paragraph{\textbf{推论 3.15.}}
\SGARefTarget{sga3:xiv:3.15}
{\itshape 设 \(G\) 与 \(H\) 如 \SGARefLink{sga3:xiv:3.14}{3.14}, 并假定对每个 \(s\in S\), 几何纤维 \(\mathfrak h_{\bar s}\) 包含 \(\mathfrak g_{\bar s}\) 的一个 Cartan 子代数. 还假定 \(\mathfrak g\) 对 fpqc 拓扑局部容许 Cartan 子代数. 则平展局部地, \(H\) 包含 \(G\) 的一个 \((\mathrm C)\) 型子群.}

由 \SGARefLink{sga3:xiv:3.7}{3.7} 和 \SGARefLink{sga3:xiv:3.9}{3.9}\textup{(a)}, \(G\) 平展局部容许 \((\mathrm C)\) 型子群. 因而可假定 \(G\) 容许一个这样的子群 \(D\). 根据共轭定理 \SGARefLink{sga3:xiii:theorem:6:1}{XIII~6.1}\textup{(a)}, 关于 \(\mathfrak h\) 的假设还说明 \SGARefLink{sga3:xiv:3.14}{3.14} 中搬运子的结构态射为满射. 结论来自 Hensel 引理 \SGARefLink{sga3:xi:proposition:1:10}{XI~1.10}.

% END inputs/XIV/tex/components/13_expose_XIV_proposition312_through_corollary315_en.tex

% BEGIN inputs/XIV/tex/components/14_expose_XIV_corollary316_remarks317_en.tex
% Authority: Expo14.pdf, local pp. 15--16 / combined pp. 849--850.
% Continuous English ends immediately before Theorem 3.18.

\paragraph{\textbf{推论 3.16.}}
\SGARefTarget{sga3:xiv:3.16}
{\itshape 设 \(G,H,K\) 如 \SGARefLink{sga3:xiv:3.12}{3.12}\textup{(a)}, 另假定每个几何纤维 \(\mathfrak k_{\bar s}\) 幂零, 即 \(\mathfrak k\) 局部幂零. 则
\[
  \operatorname{Transp}_G(K,H)
  =
  \operatorname{Transp}_G(\mathfrak k,\mathfrak h),
\]
此函子可由 \(G\) 的一个在 \(S\) 上光滑的闭子预概形表示, 其到 \(S\) 的结构态射为满射. 平展局部地, \(H\) 包含一个与 \(K\) 共轭的子群.}

两个搬运子相同仍然包含在 \SGARefLink{sga3:xiv:3.12}{3.12}\textup{(a)} 中. 关于其结构的断言恰好是 \SGARefLink{sga3:xiv:3.2}{3.2}\textup{(a)}, 而 \SGARefLink{sga3:xiv:3.16}{3.16} 的最后一个断言随即由 Hensel 引理得到.

\paragraph{\textbf{注 3.17.}}
\SGARefTarget{sga3:xiv:3.17}
\textup{a)}
在 \SGARefLink{sga3:xiv:3.12}{3.12} 和 \SGARefLink{sga3:xiv:3.16}{3.16} 中, \(K\) 在 \(S\) 上光滑的假设可以减弱为, 沿单位截面的相对 \(1\)-微分层
\[
  e_K^*(\Omega^1_{K/S})=\omega^1_K
\]
局部自由. 以这种形式, \SGARefLink{sga3:xiv:3.12}{3.12} 包含 \SGARefLink{sga3:xiii:corollary:5:5}{XIII~5.5}.

\smallskip
\noindent
\textup{b)}
设 \(G,\mathfrak g,\mathfrak h\) 如 \SGARefLink{sga3:xiv:3.1}{3.1}, 其中 \(G\) 在 \(S\) 上为有限表示. 则
\[
  N=\operatorname{Norm}_G(\mathfrak h)
\]
沿单位截面未必在 \(S\) 上光滑. 等价地, 即使 \(S\) 是域的谱, 也未必存在一个在 \(S\) 上光滑且 Lie 代数为 \(\mathfrak h\) 的子群预概形 \(H\subset G\). 当这样的 \(H\) 存在时, 取 \(H\) 具有连通纤维, 得
\[
  N=\operatorname{Norm}_G(H);
\]
我不知道此时 \(N\) 是否在 \(S\) 上光滑. 对此问题, 显然可归结到 \(S\) 为局部 Artin 概形的情形.

% END inputs/XIV/tex/components/14_expose_XIV_corollary316_remarks317_en.tex

% BEGIN inputs/XIV/tex/components/15_expose_XIV_theorem318_corollary319_en.tex
% Authority: Expo14.pdf, local pp. 16--17 / combined pp. 850--851.
% Continuous English ends immediately before Corollary 3.20.

作为本节的结束, 考察 \(G\) 在 \(S\) 上“半单”的情形.

\paragraph{\textbf{定理 3.18.}}
\SGARefTarget{sga3:xiv:3.18}
{\itshape 设 \(S\) 是预概形, \(G\) 是光滑 \(S\)-群预概形, 其几何纤维为半单“伴随型”群, 即约化中心\textup{（\SGARefLink{sga3:xii:4.1}{XII~4.1} 和 \SGARefLink{sga3:xiv:4.4}{4.4}）}为单位群的半单群. 则 \(G\) 的 \((\mathrm C)\) 型子群恰好就是它的极大环面, 因而也恰好就是它的 Cartan 子群\textup{（\SGARefLink{sga3:xii:3.1}{XII~3.1}）}.}

根据定义, 只须处理 \(S\) 是代数闭域的谱的情形. 此时须表明, 对 \(G\) 的一个极大环面 \(T\), 其 Lie 代数 \(\mathfrak t\) 是 \(\mathfrak g\) 的 Cartan 子代数. 置 \(r=\dim T=\operatorname{rank}\mathfrak t\). 考虑到不等式
\[
  \operatorname{nil\,rank}(\mathfrak g)
  \geq
  \operatorname{nil\,rank}(G)
  =
  \dim T
  =
  \operatorname{rank}\mathfrak t
  =
  r,
\]
只须找到 \(x\in\mathfrak t\), 使得
\[
  \dim\Nil(x,\mathfrak g)=r.
\]
由于 \(\mathfrak t\) 为阿贝尔 Lie 代数, 因而幂零, 这等价于寻找 \(a\in\mathfrak t\), 使得
\[
  \operatorname{ad}(a)_{\mathfrak g/\mathfrak t}
\]
为单射\textup{（\SGARefLink{sga3:xiii:corollary:5:7}{XIII~5.7}\textup{(a)}）}.

考虑由 \(G\) 的伴随表示诱导的 \(T\) 的表示中出现的 \(T\) 的特征标 \(\alpha\). 半单群 \(G\) 的结构理论\textup{（\emph{Bible}, 13, \SGARefLink{sga3:xiv:7:lemma1}{定理~1}\textup{(a)} 以及定理~3, \SGARefLink{sga3:xiv:7:lemma2}{推论~2}）}, 更确切地说 \(G\) 的“大胞腔”, 表明 \(\mathfrak g\) 中平凡特征标的特征子空间恰好是 \(\mathfrak t\). 其他特征子空间均为一维, 它们的特征标恰好是 \(G\) 相对于 \(T\) 的根. 这里的大胞腔由 \(T\) 和子群 \(P_\alpha\) 构成, 每个都同构于加法群 \(\mathbf G_a\), 在 \(T\) 下稳定, 且对应于根特征标.

\(G\) 的约化中心是 \(T\) 在伴随表示中出现的各特征标的核之交\textup{（\SGARefLink{sga3:xii:4.8}{XII~4.8}）}. 因此, \(G\) 为伴随型就是说它的根生成格
\[
  M=\operatorname{Hom}(T,\mathbf G_m).
\]
根理论中的一个熟知引理说明, 每个根都属于某个单根系, 因而属于由根生成的群的某个基, 从而属于 \(T\) 的对偶格的某个基.\footnote{重排版在此处将 \(\widehat T\) 换为 \(T\).} 得到下述结果.

\paragraph{\textbf{推论 3.19.}}
\SGARefTarget{sga3:xiv:3.19}
{\itshape 设 \(G\) 为代数闭域 \(k\) 上的伴随型半单代数群, \(T\) 为 \(G\) 的极大环面. 对 \(G\) 相对于 \(T\) 的每个根
\[
  \alpha\colon T\longrightarrow\mathbf G_m
\]
相应的同态
\[
  \alpha'\colon\mathfrak t\longrightarrow k
\]
不为零.}

此断言本质上等价于\SGARefLink{sga3:xiv:3.18}{定理 3.18}. 事实上, 对 \(t\in\mathfrak t\), 自同态 \(\operatorname{ad}(t)\) 半单, 它在 \(\mathfrak g/\mathfrak t\) 上的特征值恰好是各数 \(\alpha'(t)\). 因此 \(\operatorname{ad}(t)_{\mathfrak g/\mathfrak t}\) 为单射当且仅当所有 \(\alpha'(t)\) 均非零; 而这样的 \(t\in\mathfrak t\) 存在当且仅当每个 \(\alpha'\) 非零.

% END inputs/XIV/tex/components/15_expose_XIV_theorem318_corollary319_en.tex

% BEGIN inputs/XIV/tex/components/16_expose_XIV_corollary320_remarks321_en.tex
% Authority: Expo14.pdf, local pp. 17--19 / combined pp. 851--853.
% This closes Section 3; continuous English next begins at Section 4.

\paragraph{\textbf{推论 3.20.}}
\SGARefTarget{sga3:xiv:3.20}
{\itshape 设 \(S\) 是预概形, \(G\) 是在 \(S\) 上有限表示的光滑 \(S\)-群预概形, 其几何纤维为连通约化代数群, 即半单群由环面的扩张. 对每个 \(s\in S\), 存在 \(s\) 的一个开邻域 \(U\), 使得 \(G|_U\) 容许极大环面.}\footnote{这一结果及它所依赖的 \SGARefLink{sga3:xiv:2.11}{2.11}, 都可立即推广到将 \(s\) 换成 \(S\) 中包含在某个仿射开集内的有限子集的情形.}

我们将在第 XVI 讲看到, 关于 \(G\) 的假设蕴含 \(G\) 在 \(S\) 上仿射且约化秩局部常值. 因而由 \SGARefLink{sga3:xii:4.7}{XII~4.7}\textup{\SGARefLink{sga3:xii:4.7:item:c}{(c)}}, \(G\) 容许约化中心 \(Z\), 商
\[
  G'=G/Z
\]
是 \(S\) 上的光滑仿射群, 其约化中心为单位群, 而 \(G\) 的极大环面与 \(G'\) 的极大环面一一对应. 此外, \(G'\) 的几何纤维显然为连通半单群, 并且因其约化中心平凡, 根据定义它们是伴随型群.

所以只须处理 \(G\) 为半单伴随型群的情形. 由 \SGARefLink{sga3:xiv:3.18}{3.18}, 问题归结到寻找 \(s\) 的一个开邻域 \(U\) 和 \(G|_U\) 的一个 \((\mathrm C)\) 型子群, 等价地, 根据 \SGARefLink{sga3:xiv:3.9}{3.9}\textup{(a)}, 寻找 \(\mathfrak g|_U\) 的一个 Cartan 子代数. 当 \(k(s)\) 无限时这是可能的: 由 \SGARefLink{sga3:xiv:3.7}{3.7}, \(\mathfrak g\) 满足 \SGARefLink{sga3:xiv:2.9}{2.9} 的条件 \((\mathrm C_2)\), 可以应用 \SGARefLink{sga3:xiv:2.11}{2.11}\textup{(b)}.

事实上, 不假定 \(k(s)\) 无限, \SGARefLink{sga3:xiv:3.20}{3.20} 仍然成立. 由前述论证, 只须知道, 对有限域 \(k\) 上的每个伴随型半单群 \(G\), Lie 代数 \(\mathfrak g\) 均包含正则元素. Chevalley 利用 Weyl 群中 Coxeter 元素的性质证明了这一断言; 参见下文 J.-P. Serre 的附录.

\paragraph{\textbf{注 3.21.}}
\SGARefTarget{sga3:xiv:3.21}
\textup{a)}
若将 \(G\) 换成一个在 \(S\) 上光滑且各处的秩与 \(G\) 相同的闭子群预概形 \(H\)\textup{（例如 \(G\) 的一个“抛物子群”）}, 则只要 \(k(s)\) 无限, 断言 \SGARefLink{sga3:xiv:3.20}{3.20} 仍然成立, 证明本质上相同. 我不知道无限性的假设是否仍为多余. 至少当 \(H\) 为抛物子群时可以去掉它: 构造 \(H\) 的根基 \(U\) 及商 \(H/U\), 就将问题归结到半单情形. 遗憾的是, 正则元素的方法在此似乎无能为力. 容易构造出一些例子——例如射影群及其“标准”Borel 子群——使得 \(H\) 的 Lie 代数 \(\mathfrak h\) 不包含任何正则元素.

\smallskip
\noindent
\textup{b)}
当 \(k(s)\) 无限时, \SGARefLink{sga3:xiv:3.20}{3.20} 的证明结合 \SGARefLink{sga3:xiv:3.9}{3.9}\textup{(b)} 实际给出如下更强的陈述: \(G_s\) 的每个极大环面 \(T_0\) 均由 \(s\) 的某个开邻域上的极大环面 \(T\) 诱导. 我不知道去掉无限性的假设后这是否仍然成立. 显然的困难在于, \(T_0\) 的 Lie 代数 \(\mathfrak t_0\) 未必包含纤维 \(G_s\) 的 Lie 代数 \(\mathfrak g_0\) 的正则元素.

肯定的答案将蕴含下述断言, 目前只在剩余域无限或 \(A\) 分离且完备时已知. 设 \(A\) 为局部环, 剩余域为 \(k\); 设 \(M\) 为 \(A\) 上的一个“Azumaya 代数”, 即作为 \(A\)-模为有限自由模, 且其约化
\[
  M_0=M\otimes_A k
\]
是 \(k\) 上的中心单代数的代数. 设 \(D_0\subset M_0\) 为一个在 \(k\) 上可分的交换子代数, 满足
\[
  [M_0:k]=[D_0:k]^2.
\]
则应存在一个作为 \(A\)-模为直和项的交换子代数 \(D\subset M\), 满足
\[
  D\otimes_A k=D_0.
\]

给定 \(M\) 等价于给定一个射影群 \(\operatorname{PGL}(n)_A\) 下的主齐性丛, 从而等价于给定 \(\operatorname{PGL}(n)\) 的一个内扭曲形式 \(G\). \(G\) 的极大环面与在 \(A\) 上平展且秩为 \(n\) 的交换子代数 \(D\subset M\) 一一对应.

\smallskip
\noindent
\textup{c)}
将 \SGARefLink{sga3:xiv:3.20}{3.20} 应用于约化 \(S\)-群预概形 \(G\) 的子环面 \(Q\) 的中心化子, 可见每个这样的 \(Q\) 均 Zariski 局部包含在 \(G\) 的一个极大环面中.

% END inputs/XIV/tex/components/16_expose_XIV_corollary320_remarks321_en.tex

% BEGIN inputs/XIV/tex/components/17_expose_XIV_section4_definition41_through_definition45_en.tex
% Authority: Expo14.pdf, local pp. 18--19 / combined pp. 852--853.
% Exact next cursor: Remarks 4.6 on local p. 19 / combined p. 853.

\section{关于 Borel 子群的插论}
\SGARefTarget{sga3:xiv:4}

\paragraph{\textbf{定义 4.1.}}
\SGARefTarget{sga3:xiv:4.1}
{\itshape 设 \(G\) 是代数闭域上的光滑代数群. \(G\) 的 \emph{Borel 子群}是具有光滑、可解、连通这些性质且对这些性质极大的代数子群.}

当 \(G\) 仿射时, 这与\textup{\emph{(Bible)}, 6, \SGARefLink{sga3:xiv:7:lemma1}{定义~1}} 中的术语一致. 立即注意到, 若 \(Z\) 为 \(G\) 的一个包含在其中心内的连通光滑子群\textup{（或更一般地, 若 \(Z\) 可解且正规）}, 则对 \(G\) 的每个 Borel 子群 \(B\), \(B\times Z\) 在从 \(B\times Z\) 到 \(G\) 的态射
\[
  (b,z)\longmapsto bz
\]
下的像 \(BZ\) 是 \(G\) 的一个包含 \(B\) 的光滑、可解、连通子群. 因而它等于 \(B\). 所以 \(B\) 包含 \(Z\), 从而 \(B\) 是 \(G'=G/Z\) 的某个代数子群 \(B'\) 的逆像; 立即可知 \(B'\) 为 \(G'\) 的 Borel 子群.

取
\[
  Z=\Centr(G^0)_{\mathrm{red}}^0,
\]
由 \SGARefLink{sga3:xii:6.1}{XII~6.1}, 商 \(G'=G/Z\) 仿射. 因此 \(G'\) 的 Borel 子群彼此共轭, 且对每个这样的 \(B'\), \(G'/B'\) 均为射影簇\textup{\emph{(Bible)}, 6, \SGARefLink{sga3:xiv:7:lemma4}{定理~4}}. 从而:

\paragraph{\textbf{命题 4.2.}}
\SGARefTarget{sga3:xiv:4.2}
{\itshape 设 \(G\) 如 \SGARefLink{sga3:xiv:4.1}{4.1}. 则 \(G\) 的 Borel 子群彼此共轭. 若 \(B\) 是 Borel 子群, 则 \(G/B\) 为射影簇. \(B\) 的极大环面是 \(G\) 的极大环面; 当 \(G\) 连通时, \(B\) 的 Cartan 子群是 \(G\) 的 Cartan 子群.}

只须证明最后一个断言, 并且显然可假定 \(G=G^0\). 关于 Cartan 子群的断言, 来自 \(G'\) 中的类似断言\textup{\emph{(Bible)}, 6, \SGARefLink{sga3:xiv:7:lemma4}{定理~4}, \SGARefLink{sga3:xiv:7:lemma4}{推论~4}} 以及 \SGARefLink{sga3:xii:6.6}{XII~6.6}\textup{\SGARefLink{sga3:xii:6.6:item:e}{(e)}}. 关于极大环面的断言由此推出, 因为根据 \SGARefLink{sga3:xii:6.6}{XII~6.6}\textup{\SGARefLink{sga3:xii:6.6:item:c}{(c)}}, 光滑代数群的极大环面恰好就是其 Cartan 子群的极大环面.

\paragraph{\textbf{推论 4.3.}}
\SGARefTarget{sga3:xiv:4.3}
{\itshape 假定 \(G\) 连通. 则 \(G\) 的每个元素都包含在 \(G\) 的某个 Borel 子群中.}

这归结到 \(G'\) 中的同一断言, 后者是熟知的\textup{\emph{(Bible)}, 6, 定理~5\textup{(d)}}.

\paragraph{\textbf{推论 4.4.}}
\SGARefTarget{sga3:xiv:4.4}
{\itshape 设 \(B\) 是 \(G\) 的 Borel 子群, \(C\) 是 \(B\) 的 Cartan 子群, \(N\) 是它在 \(G\) 中的正规化子. 则
\[
  N\cap B=C.
\]}

事实上,
\[
  N\cap B=\operatorname{Norm}_B(C),
\]
所以只须表明, 当 \(G\) 连通且可解时, Cartan 子群 \(C\) 等于自身的连通正规化子. 沿用上述记号, \(C\) 是 \(G'\) 的某个 Cartan 子群 \(C'\) 的逆像, 因而归结到 \(G\) 仿射的情形. Cartan 子群的正规化子光滑\textup{（\(C\) 为其自身的连通正规化子; 例如参见 \SGARefLink{sga3:xii:6.6}{XII~6.6}\textup{\SGARefLink{sga3:xii:6.6:item:c}{(c)}}）}. 所以只须检验 \(C\) 与 \(N\) 具有相同的 \(k\)-值点, 这恰好就是\textup{\emph{(Bible)}, 定理~6\textup{\SGARefLink{sga3:xii:6.6:item:d}{(d)}}}.

\paragraph{\textbf{定义 4.5.}}
\SGARefTarget{sga3:xiv:4.5}
{\itshape 设 \(G\) 是预概形 \(S\) 上的有限表示光滑群预概形. \(G\) 的一个 \emph{Borel 子群}是 \(G\) 的有限表示光滑子群预概形 \(B\), 使得对每个 \(s\in S\), 几何纤维 \(B_{\bar s}\) 是 \(G_{\bar s}\) 的 Borel 子群.}

立即可验证, 这一概念在基变换下稳定, 并对 fpqc 拓扑是局部的. 事实上, 若 \(k'\) 是代数闭域 \(k\) 的代数闭扩张, 则 \(k\) 上光滑代数群 \(G\) 的代数子群 \(B\) 是 \(G\) 的 Borel 子群当且仅当 \(B_{k'}\) 是 \(G_{k'}\) 的 Borel 子群.

由定义可知, 若 \(G\) 是任意域 \(k\) 上的光滑代数群, \(B\) 是 \(G\) 的 Borel 子群, 则 \(G/B\) 为射影簇; \(B\) 的每个极大环面 \(T\) 都是 \(G\) 的极大环面; \(T\) 在 \(B\) 中的正规化子等于其中心化子 \(C\), 当 \(G\) 连通时, 后者是 \(G\) 的 Cartan 子群.

% END inputs/XIV/tex/components/17_expose_XIV_section4_definition41_through_definition45_en.tex

% BEGIN inputs/XIV/tex/components/18_expose_XIV_remarks46_through_remarks481_en.tex
% Authority: Expo14.pdf, local pp. 19--20 / combined pp. 853--854.
% The source has no item numbered 4.7.
% Exact next cursor: the paragraph preceding Proposition 4.9,
% local p. 20 / combined p. 854.

\paragraph{\textbf{注 4.6.}}
\SGARefTarget{sga3:xiv:4.6}
遗憾的是, 一般而言, \(G\) 的两个 Borel 子群对 fpqc 拓扑局部共轭不再成立. 即使 \(G\) 在 \(S\) 上仿射且 \(S\) 是代数闭域 \(k\) 上对偶数代数的谱, 这一点也可能失败.

作为这一遗憾事实的一个后果, 应指出, 若 \(G\) 是非完美域 \(k\) 上的光滑仿射连通代数群, 则一般不可能自然地定义一个 \(G\) 下的齐性空间 \(D\), 以起到旗簇, 即 \(G\) 的 Borel 子群簇的作用. 在 \(k\) 的代数闭包 \(\bar k\) 上, 这样的空间须同构于 \(G_{\bar k}/B\), 其中 \(B\) 是 \(G_{\bar k}\) 的 Borel 子群.

事实上, 假定 \(G_{\bar k}\) 对其根基 \(\bar R\) 的商为伴随型半单群. 则
\[
  G_{\bar k}\longrightarrow
  \operatorname{Aut}_{\bar k}(\bar D)
\]
的核为 \(G_{\bar k}\) 的根基. 因此, 若 \(\bar D\) 来自一个 \(G\) 下的齐性空间 \(D\), 根基 \(\bar R\) 就来自 \(G\) 的一个子群 \(R\). 但容易构造出一些例子, 使 \(G_{\bar k}/\bar R\) 为伴随型, 而 \(\bar R\) 并未定义在 \(k\) 上.

在这些条件下, 函子
\[
  \mathfrak B\colon(\operatorname{Sch}/S)^\circ
    \longrightarrow \operatorname{Ens},
  \qquad
  \mathfrak B(S')=
  \{G_{S'}\text{ 的 Borel 子群}\},
\]
不能由一个光滑 \(S\)-预概形表示. 从无穷小观点\textup{（III, \S3）}来看, 共轭定理的失败表现为如下事实: 若 \(B\) 为光滑代数群 \(G\) 的 Borel 子群, 则上同调群
\[
  H^1(B,\mathfrak g/\mathfrak b)
\]
可能不为零.\footnote{这里 \(\mathfrak b\) 表示 \(B\) 的 Lie 代数.}

另一方面, 我们将在后面的一讲看到, 当 \(G\) 半单, 或更一般地约化时, 这些不理想的现象不会出现. 这些现象, 加上无疑同样相关的缺乏良好存在性定理这一点, 解释了为何在一般群概形的概形论研究中, Borel 子群只起到相对不显著的作用, 而在后面各讲的半单群概形理论中却占据主导地位.

\paragraph{\textbf{命题 4.8.}}
\SGARefTarget{sga3:xiv:4.8}
\footnote{底本中没有编号为 4.7 的条目.}
{\itshape 设 \(G\) 是具有连通纤维的有限表示光滑 \(S\)-群预概形, \(B\) 是 \(G\) 的 Borel 子群. 则 \(B\) 等于自身的正规化子, 并且是 \(G\) 的闭子预概形.}

由 \SGARefLink{sga3:xii:7.10}{XII~7.10}, 只须证明, 在代数闭域 \(k\) 上, \(G(k)\) 中每个正规化 \(B\) 的元素均属于 \(B(k)\). 当 \(G\) 仿射时, 这是 Chevalley 的一个基本结果\textup{\emph{(Bible)}, 9, \SGARefLink{sga3:xiv:7:lemma1}{定理~1}}; 一般情形利用 \SGARefLink{sga3:xiv:4.2}{4.2} 中已用过的归约立即归结于此.

\paragraph{\textbf{注 4.8.1.}}
\SGARefTarget{sga3:xiv:4.8.1}
还可以引入 \(G\) 的\emph{抛物子群}这一概念, 以推广\SGARefLink{sga3:xiv:4.5}{定义~4.5}. 它是 \(G\) 的一个在 \(S\) 上光滑且有限表示的子群预概形 \(P\), 满足对每个 \(s\in S\), 几何纤维 \(P_{\bar s}\) 是 \(G_{\bar s}\) 的抛物子群, 即包含 \(G_{\bar s}\) 的一个 Borel 子群.

\SGARefLink{sga3:xiv:4.8}{命题~4.8} 可推广到 \(G\) 的抛物子群 \(P\), 仍然归结到已知的集合论断言. 记下如下后果\textup{（比较第 XVI 讲）}. \SGAInlineRefTarget{sga3:xiv:paragraph:4-8-bis}若 \(P\) 为 \(G\) 的抛物子群, 则 \(G/P\) 可由 \(S\) 上有限表示的拟射影预概形表示; 此处假定 \(G\) 具有连通纤维. 此外, \(G/P\) 显然在 \(S\) 上光滑, 且其几何纤维连通且固有. 容易由 EGA~III, 5.5.1 推出
\[
  D=G/P
\]
在 \(S\) 上固有, 因而射影. 若其相对维数为 \(n\), 已知可逆层 \(\Omega^n_{D/S}\) 的逆层在 \(D/S\) 的几何纤维上诱导丰沛层; 因而由 EGA~III, 4.7.1,
\[
  \bigl(\Omega^n_{D/S}\bigr)^{-1}
\]
在 \(D\) 上相对于 \(S\) 丰沛.

% END inputs/XIV/tex/components/18_expose_XIV_remarks46_through_remarks481_en.tex

% BEGIN inputs/XIV/tex/components/19_expose_XIV_proposition49_corollary410_en.tex
% Authority: Expo14.pdf, local pp. 20--21 / combined pp. 854--855.
% Exact next cursor: Theorem 4.11, local p. 21 / combined p. 855.

归结到仿射情形及代数闭基域后容易看出, 若
\[
  u\colon G\longrightarrow G'
\]
是光滑代数群之间的满态射, 则对 \(G\) 的每个 Borel 子群 \(B\), \(u(B)=B'\) 均为 \(G'\) 的 Borel 子群. 现在考察这一对应在 \(G\) 的 Borel 子群与 \(G'\) 的 Borel 子群之间给出双射的情形.

\paragraph{\textbf{命题 4.9.}}
\SGARefTarget{sga3:xiv:4.9}
{\itshape 设 \(G\) 与 \(G'\) 为具有连通纤维的有限表示光滑 \(S\)-群预概形, 设
\[
  u\colon G\longrightarrow G'
\]
为忠实平坦, 即满射的群同态.\footnote{编者询问此处是否应补充一个引用.} 置 \(N=\ker u\), 并假定下列两个条件之一成立:

\smallskip
\noindent
\textup{a)} \(N\) 包含在 \(G\) 的中心内.

\smallskip
\noindent
\textup{b)} 对某个域 \(k\), 有 \(S=\operatorname{Spec}(k)\); 若 \(\bar k\) 为 \(k\) 的代数闭包, 则
\[
  N_{\bar k}=\ker(u_{\bar k})
\]
包含在 \(G_{\bar k}\) 的根基中, 即包含在其最大的光滑、连通、可解、正规子群中.

\smallskip
则映射
\[
  B'\longmapsto u^{-1}(B')
\]
诱导从 \(G'\) 的 Borel 子群集合到 \(G\) 的 Borel 子群集合的双射.}

情形 b) 立即来自 \(G'\) 的代数子群与 \(G\) 中包含 \(N\) 的代数子群之间的对应, 以及在代数闭域上 \(G\) 的 Borel 子群包含其根基这一事实. 后一事实由 \SGARefLink{sga3:xiv:4.2}{4.2} 之前的论证立即得到.

证明情形 a). 由 \SGARefLink{sga3:xii:7.12}{XII~7.12}, 映射
\[
  H'\longmapsto H=u^{-1}(H')
\]
在 \(G'\) 的如下子群预概形 \(H'\) 与 \(G\) 的具有类似性质的子群预概形 \(H\) 之间给出双射: 它们在 \(S\) 上光滑且有限表示, 具有连通纤维, 并在每个 \(s\in S\) 处与 \(G'\) 的约化秩、幂零秩均相同. \(G'\) 和 \(G\) 的 Borel 子群具有这些性质.

只须表明, 当 \(H'\) 与 \(H\) 相对应时, \(H'\) 为 \(G'\) 的 Borel 子群当且仅当 \(H\) 为 \(G\) 的 Borel 子群. 由定义, 此问题归结到 \(S\) 是代数闭域的谱的情形. 由于 \(N\) 包含在 \(G\) 的中心内, 因而也包含在 \(H\) 的中心内, 群 \(H\) 可解当且仅当 \(H'\) 可解. 最后, \(G'\) 的代数子群与 \(G\) 中包含 \(N\) 的代数子群之间的对应立即表明, \(H'\) 满足\SGARefLink{sga3:xiv:4.1}{定义~4.1} 的极大性当且仅当 \(H\) 满足此性质. 证明完成.

\paragraph{\textbf{推论 4.10.}}
\SGARefTarget{sga3:xiv:4.10}
{\itshape 沿用 \SGARefLink{sga3:xiv:4.9}{4.9} 的记号,\footnote{底本印为“\SGARefLink{sga3:xiv:4.9}{4.7}”; 所指应为 \SGARefLink{sga3:xiv:4.9}{4.9}.} 若 \(B'\) 与 \(B\) 为 \(G'\) 与 \(G\) 中相对应的 Borel 子群, 则
\[
  \mathfrak b=\operatorname{Lie}(u)^{-1}(\mathfrak b'),
\]
其中 \(\mathfrak g,\mathfrak g',\mathfrak b,\mathfrak b'\) 分别为 \(G,G',B,B'\) 的 Lie 代数, 而
\[
  \operatorname{Lie}(u)\colon\mathfrak g\longrightarrow\mathfrak g'
\]
为 \(u\) 诱导的同态.}

这由定义及恒等式
\[
  u^{-1}(B')=B.
\]

立即得到. 现在可以证明本节的主要结果.

% END inputs/XIV/tex/components/19_expose_XIV_proposition49_corollary410_en.tex

% BEGIN inputs/XIV/tex/components/20_expose_XIV_theorem411_corollary412_en.tex
% Authority: Expo14.pdf, local pp. 21--22 / combined pp. 855--856.
% This closes Section 4.
% Exact next cursor: Section 5 and Theorem 5.1,
% local p. 22 / combined p. 856.

\paragraph{\textbf{定理 4.11.}}
\SGARefTarget{sga3:xiv:4.11}
{\itshape 设 \(G\) 是代数闭域 \(k\) 上的光滑代数群, \(\mathfrak g\) 为其 Lie 代数. 则 \(\mathfrak g\) 是 \(G\) 的各 Borel 子群 \(B\) 的 Lie 代数 \(\mathfrak b\) 的并.}

显然可假定 \(G\) 连通. 令 \(R\) 为 \(G\) 的根基, 并置 \(G'=G/R\). 由 \SGARefLink{sga3:xiv:4.9}{4.9}\textup{(b)} 和 \SGARefLink{sga3:xiv:4.10}{4.10}, 只须对 \(G'\) 证明该定理; 所以可假定 \(G\) 半单.

令 \(Z\) 为 \(G\) 的中心, 它就是其约化中心, 并置 \(G'=G/Z\). 同一论证, 这次利用 \SGARefLink{sga3:xiv:4.9}{4.9}\textup{(a)}, 将定理归结到 \(G'\). 因而可假定 \(G\) 为半单伴随型群.

令 \(B\) 为 \(G\) 的 Borel 子群, \(T\) 为 \(B\) 的一个极大环面, 因而也是 \(G\) 的极大环面; 令 \(\mathfrak b\) 与 \(\mathfrak t\) 为它们的 Lie 代数. 由 \SGARefLink{sga3:xiv:3.18}{3.18}, \(T\) 是 \(G\) 的一个 \((\mathrm C)\) 型子群, 即 \(\mathfrak t\) 是 \(\mathfrak g\) 的 Cartan 子代数. 所以 \(\mathfrak t\) 的各共轭子代数之并在 \(\mathfrak g\) 中稠密\textup{（\SGARefLink{sga3:xiii:theorem:5:1}{XIII~5.1}, \textup{(i)}\(\Rightarrow\)\textup{(vii)}）}. \(\mathfrak b\) 的各共轭子代数之并更在 \(\mathfrak g\) 中稠密.

令 \(X\) 为
\[
  G/B\times W(\mathfrak g)
\]
的闭子概形, 它的 \(k\)-值点为满足下式的二元组 \((gB,x)\):
\[
  x\in\operatorname{Ad}(g)\mathfrak b
\]
\textup{（比较 XIII~1）}. 由第二投影诱导的态射
\[
  \psi\colon X\longrightarrow W(\mathfrak g)
\]
固有, 因为 \(G/B\) 在 \(k\) 上固有. 刚才已看到它占优; 因而它为满射. 这就证明了 \SGARefLink{sga3:xiv:4.11}{4.11}.

在本讲余下部分中, 将用到的本节结果只有下述推论.

\paragraph{\textbf{推论 4.12.}}
\SGARefTarget{sga3:xiv:4.12}
{\itshape 设 \(k\) 为无限域, \(G\) 为 \(k\) 上的光滑代数群, \(T\) 为 \(G\) 的极大环面. 以 \(\mathfrak g\) 与 \(\mathfrak t\) 表示它们的 Lie 代数. 设
\[
  u\colon G\longrightarrow\operatorname{GL}(V)
\]
为线性表示, 其中 \(V\) 在 \(k\) 上有限维, 设
\[
  u'\colon\mathfrak g\longrightarrow\mathfrak{gl}(V)
\]
为所诱导的 Lie 代数表示. 则当 \(x\in\mathfrak g\) 变化时, \(u'(x)\) 的零度的最小值可在某个元素 \(x\in\mathfrak t\) 处达到.}

立即归结到 \(k\) 代数闭的情形. 显然可假定 \(G\) 连通, 并且将 \(G\) 对 \((\ker u)_{\mathrm{red}}\) 取商后, 可假定 \(G\) 仿射. 利用 \SGARefLink{sga3:xiv:4.11}{4.11} 和 \(G\) 的极大环面的共轭定理, 进一步归结到 \(G\) 可解的情形.

此时 \(G\) 为半直积 \(T\cdot U\), 其中 \(U\) 为 \(G\) 的幺幂部分, 是一个光滑连通幺幂群\textup{\emph{(Bible)}, 6, 定理~3}. 因此, 作为向量空间,
\[
  \mathfrak g=\mathfrak t\oplus\mathfrak u,
  \qquad
  \mathfrak u=\operatorname{Lie}(U).
\]
由 Lie--Kolchin 定理\textup{\emph{(Bible)}, 6, \SGARefLink{sga3:xiv:7:lemma1}{定理~1}}, 表示空间 \(V\) 具有由稳定子空间 \(V_i\) 组成的合成列, 其商
\[
  W_i=V_i/V_{i+1}.
\]
均为一维. 因而对每个 \(i\), 存在诱导的表示
\[
  u_i\colon G\longrightarrow
  \operatorname{GL}(W_i)=\mathbb G_m
\]
以及相应的 Lie 代数同态
\[
  u_i'\colon\mathfrak g\longrightarrow k.
\]
对每个 \(x\in\mathfrak g\), \(u'(x)\) 的零度等于满足 \(u_i'(x)=0\) 的指标 \(i\) 的个数.

由于 \(U\) 幺幂, 每个 \(u_i\) 在 \(U\) 上均平凡, 因而每个 \(u_i'\) 在 \(\mathfrak u\) 上为零. 若
\[
  x=t+v,\qquad t\in\mathfrak t,\quad v\in\mathfrak u,
\]
则对每个 \(i\), 均有 \(u_i'(x)=u_i'(t)\). 因此 \(u'(x)\) 与 \(u'(t)\) 的零度相同, \SGARefLink{sga3:xiv:4.12}{4.12} 得证.

% END inputs/XIV/tex/components/20_expose_XIV_theorem411_corollary412_en.tex
